\documentclass[amsmath,amsfonts,amssymb,aps,pra,superscriptaddress,10pt]{revtex4-2}
\usepackage[margin=25mm]{geometry}
\usepackage[T1]{fontenc}
\usepackage{lmodern}
\usepackage{amsmath,amssymb,amsthm,mathtools}
\usepackage{mathrsfs}
\usepackage{caption}
\captionsetup{compatibility=false}
\usepackage{xcolor,microtype}
\usepackage[labelsep=period]{caption}
\usepackage[colorlinks=true,allcolors=blue!45!black]{hyperref}
\hypersetup{pdfauthor={Davide Batic and Denys Dutykh}}
\setlength{\emergencystretch}{2em}
\setlength{\parskip}{2pt}
\allowdisplaybreaks[1]
\numberwithin{equation}{section}
\theoremstyle{plain}
\newtheorem{theorem}{Theorem}
\newtheorem{proposition}{Proposition}
\newtheorem{lemma}{Lemma}
\newtheorem{corollary}{Corollary}
\theoremstyle{definition}
\newtheorem{definition}{Definition}
\newtheorem{background}{Background}[section]
\theoremstyle{remark}
\newtheorem{remark}{Remark}
\newcommand{\ellzero}{\ell_0}
\newcommand{\RG}{R_{\mathrm G}}
\newcommand{\vc}{v_{\mathrm c}}
\newcommand{\Drz}{\Delta_{rz}}
\newcommand{\gradRZ}{\nabla_{rz}}
\newcommand{\ET}{\mathcal T_{\ellzero}}
\newcommand{\EE}{\mathcal E_{\ellzero}}
\newcommand{\OmZ}{\Omega_{\mathrm Z}}

\begin{document}

\title{Local realizability and global completion of constant-\texorpdfstring{$\ell$}{ell} rotating dust}
\author{Davide Batic}
\email{davide.batic@ku.ac.ae}
\affiliation{Mathematics Department, Khalifa University of Science and Technology, PO Box 127788, Abu Dhabi, United Arab Emirates}

\author{Denys Dutykh}
\email{denys.dutykh@ku.ac.ae}
\affiliation{Mathematics Department, Khalifa University of Science and Technology, PO Box 127788, Abu Dhabi, United Arab Emirates}

\date{\today}

\begin{abstract}
We study local realizability and global completion for stationary, axisymmetric rotating dust in the class characterised by a constant logarithmic derivative of the metric potential. An exact point transformation reduces the velocity equation to a linear elliptic equation. The inverse transformation develops a fold at the lower density boundary, while positive density selects a regular sheet. Flat equatorial rotation curves admit local analytic spacetime realisations on annuli. Their continuation is constrained by a condition on metric derivatives at the axis and by quantitative radial-span bounds under an invariant transverse-concavity hypothesis. These bounds are sharp for nonpositive values of the defining constant and, in the exactly flat limit, for positive values. For a nonzero defining constant, a uniformly subluminal branch on a standard asymptotically flat end has inverse-radius velocity decay with a positive leading coefficient. Nonnegative density near the equatorial tail then forces a locally flat vacuum region. For finite bodies, toroidal topology imposes period conditions on exterior potentials, and completion without a surface layer requires the full Darmois matching data. We derive a compatibility condition for the gauge-fixed linearised Ernst equations and identify the additional data needed for spectroscopic prediction. Together, these results clarify the conditions governing the passage from local rotation profiles to a complete isolated spacetime.
\end{abstract}

\maketitle

\section{Introduction}\label{Intro}

Nearly flat orbital-speed profiles in the outer regions of disk galaxies are a central feature of the galactic mass discrepancy \cite{Lelli2016AJ}. Their standard interpretation invokes a substantial non-luminous matter component \cite{BertoneHooper2018RMP}. Relativistic models of rotating matter provide a setting in which to examine how the inferred mass distribution depends on spacetime geometry, differential rotation, and the choice of observer. Such an investigation requires a connection between local kinematics and a complete gravitational configuration. A prescribed speed profile must satisfy the field equations, arise from a physically admissible matter distribution, and extend consistently to the axis or to an exterior region. Comparison with observations further requires the frequency shifts associated with specified emitters, observers, and connecting null geodesics. These constraints make the relation between local realizability and global completion particularly important for pressureless rotating dust. The stationary, axisymmetric dust family described by Winicour can be parametrized by a scalar field $\eta(r,z)$ and a negative function $H(\eta)$ \cite{Winicour1975JMP, Stephani2003Exact}. The speed measured by a zero-angular-momentum observer (ZAMO) is $v=\eta/r$, while differential rotation is encoded by $\ell=H_{,\eta}/H$. Balasin and Grumiller used the rigidly rotating sector of this family in a galactic model \cite{Balasin2008IJMPD}. Astesiano, Cacciatori, Gorini, and Re subsequently developed a differential-rotation construction and considered a constant value of $\ell$ near a galactic edge \cite{Astesiano2022EPJC}. Their analysis combines a low-velocity truncation with a near-edge replacement of the radial coefficients. The resulting equation admits a nonzero constant speed, which motivates an examination of the corresponding exact branch. In particular, the relation between this approximate construction and the parent equations depends on how the prescribed equatorial profile is extended in the transverse direction.

Global restrictions on rotating dust have a substantial history. Zingg, Aste, and Trautmann derived finite-radius density constraints for the general non-rigid Winicour family and examined the simultaneous prescription of metric and normal data at a vacuum interface \cite{ZinggAsteTrautmann2007ASTP}. Their analysis established the importance of treating matter support and exterior matching as linked parts of the same problem. Caporali's global result, Pfister's conditional analysis, and the singularity findings for rigid van Stockum--Bonnor configurations provide further precedents \cite{Caporali1978PLA, Pfister2010CQG, BratekJalochaKutschera2007PRD}. The present work builds on this foundation by deriving explicit criteria within the constant-$\ell$ subclass. The resulting density interval, reconstruction geometry, and radial-span estimates give a quantitative description of how local solutions encounter the restrictions associated with increasing radial extent. A complementary literature examines the magnitude and interpretation of relativistic effects in galactic dynamics. Calculations for physically sourced weak-field disks quantify the small gravitomagnetic corrections and emphasize the importance of consistent source modelling and approximation schemes \cite{Ciotti2022ApJ, Lasenby2023CQG, GlampedakisJones2023CQG}. Analyses of the rigidly rotating model relate its frame dragging to the asymptotic reference frame and axial structure \cite{CostaEtAl2023PRD}. Broader studies of stationary geometry and galactic observables provide further constraints \cite{CostaNatario2024PRD, Ruggiero2024JCAP}. Our analysis complements these approaches through an exact study of the differential-rotation branch. It follows the reconstruction from a prescribed local trace through density admissibility, geometric regularity, and exterior boundary data. The present analysis identifies the precise assumptions underlying each restriction on the continuation of local solutions. The central question is whether a constant-$\ell$ branch can realize an exactly or approximately flat equatorial speed profile and support its continuation to a regular isolated spacetime. An equatorial trace leaves transverse derivatives available to satisfy the field equations. This freedom is essential to the local construction. Direct substitution of a spatially constant nonzero speed produces a residual in both the exact and retained-order equations, while suitable transverse dependence permits genuine local realizations of a flat equatorial trace. Keeping track of this dependence gives a precise mathematical formulation of the passage from a rotation profile to a dust spacetime.

We begin with the exact linearizing potential and examine the inverse transformation needed to recover the physical variables. The density inequality determines the admissible range and identifies a fold at its lower boundary. Positive density selects a regular inverse sheet on which analytic equatorial data can be reconstructed into a local spacetime metric. We prove the existence of analytic, positive-density realizations of exactly flat traces on admissible annuli and establish automatic analyticity on regular sheets. We also exhibit Hadamard instability for the auxiliary elliptic Cauchy problem. Together, these results describe both the local freedom of the construction and the sensitivity of continuation from equatorial data. The continuation analysis then introduces geometric conditions that control this local freedom. Under the stated regularity and timelikeness assumptions, a metric argument at an effective rotation axis forces $H'(0)=0$. A constant-$\ell$ branch reaching such an axis therefore has $\ell=0$. For reflection-symmetric solutions on annuli, nonnegative density and an invariant transverse-concavity condition for the relative ZAMO rapidity yield quantitative bounds on the radial span of exactly or approximately flat traces. These estimates relate the permitted departure from flatness to the extent of the annulus under a specified geometric shape hypothesis. We establish sharpness for nonpositive $\ell$ and for positive $\ell$ in the exactly flat limit. This gives a quantitative refinement of the earlier finite-radius restrictions within the class considered here. Global completion leads to two further problems. For a regular branch with nonzero constant $\ell$ extending along a standard asymptotically flat equatorial end, uniform subluminality and the exact metric relations force $rv$ to approach a positive constant. Nonnegative density on open neighbourhoods of the equatorial tail then forces a locally flat vacuum region. Analytic continuation propagates the resulting constant auxiliary matter variable across a connected regular inverse sheet. The argument identifies the limiting geometry directly from the exact equations and clarifies the conditions under which an asymptotic extrapolation is compatible with the matter inequality. For a finite body, completion requires an exterior realizing the boundary geometry supplied by the interior. When the body has a disk-shaped meridional section separated from the axis, its toroidal topology introduces period conditions for globally defined exterior potentials. At the interface, matching without a surface layer requires the full Darmois data \cite{Darmois1927,Israel1966NCB,MarsSenovilla1993CQG}. A vanishing density at the boundary leaves these geometric requirements to be satisfied. We formulate an additional compatibility condition through the linearized Ernst equations about Minkowski space on a fixed smooth bounded domain away from the axis. In a specified gauge, the associated Dirichlet-to-Neumann map characterizes the admissible pairs of boundary values and conormal derivatives. This gives a concrete boundary test within the linearized problem. Construction of a nonlinear isolated exterior remains a separate existence question. Finally, the observational interpretation connects the completed geometry to the emitter--observer frequency ratio. The local ZAMO speed enters this relation together with the emitter and observer motions and the conserved quantities of the connecting null ray. Stating these dependencies identifies the additional information required to turn a locally realized rotation profile into a spectroscopic prediction. The overall picture combines a constructive local theory with explicit conditions governing continuation and measurement.

The manuscript is organized as follows. Sections~\ref{sec:model-framework}--\ref{sec:approx-flat-constraint} develop the local equations and continuation bounds. Sections~\ref{sec:asymptotic-constraint}--\ref{sec:topology-exterior} address asymptotic and finite completion. Section~\ref{sec:observable-scope} relates the resulting geometry to spectroscopic data. The supplementary material gives the Runge approximation and reconstruction details, an alternative axis argument, an explicit flat-trace solution, further radial-span estimates and sharpness constructions, comparisons with the retained-order equations, additional junction formulas, the hypotheses for nonlinear exterior matching, and a detailed formulation of the spectroscopic prediction problem. It also describes the exact algebraic checks accompanying the source package. Finally, every proof needed for the principal results stated in this work is included below.


\section{The exact constant-\texorpdfstring{$\ell$}{ell} equations}\label{sec:model-framework}
We use signature $(-,+,+,+)$, units $c=1$, and coordinates $(t,\phi,r,z)$ adapted to the stationary and axial Killing fields $k=\partial_t$ and $m=\partial_\phi$. The axial orbits have period $2\pi$. In the circular rotating-dust family considered here, the line element is
\begin{equation}\label{eq:model-block-metric}
ds^2=G_{AB}\,dx^A dx^B+e^\mu(dr^2+dz^2),\qquad
G=(g_{AB})_{A,B\in\{t,\phi\}}.
\end{equation}
The meridional coordinates use the canonical orbit-area radius
\begin{equation}\label{eq:model-metric-identities}
r=\sqrt{-\det G},\qquad \det G=-r^2.
\end{equation}
This gauge is part of the exact dust representation \cite{Winicour1975JMP,Astesiano2022EPJC}. The matter tensor is $T_{\mu\nu}=\rho u_\mu u_\nu$, with circular four-velocity
\begin{equation}\label{eq:model-velocity-relations}
u=\frac{k+\Omega m}{\sqrt{-H}},\qquad
\frac{d\Omega}{d\eta}=\frac{H_{,\eta}}{2\eta}.
\end{equation}
Moreover, we have
\begin{equation}\label{eq:model-H-eta-definitions}
H=g(k+\Omega m,k+\Omega m)<0,\qquad
\eta=g(k+\Omega m,m),
\end{equation}
and the Killing block is
\begin{equation}\label{eq:model-metric-components}
g_{\phi\phi}=\frac{r^2-\eta^2}{-H},\qquad
g_{t\phi}=\eta-\Omega g_{\phi\phi},\qquad
g_{tt}=H-2\Omega\eta+\Omega^2g_{\phi\phi}.
\end{equation}
The remaining function $\mu$ follows from the quadratures established in the local reconstruction below. For a hypersurface with spacelike unit normal $n$, our second fundamental form convention is
\begin{equation}\label{eq:extrinsic-curvature-convention}
K_{ab}=h_a{}^\mu h_b{}^\nu\nabla_\mu n_\nu.
\end{equation}
The ZAMO angular velocity and lapse satisfy
\begin{equation}\label{eq:zamo-shift-definition}
\OmZ=-\frac{g_{t\phi}}{g_{\phi\phi}},\qquad
N^2=-g_{tt}+\frac{g_{t\phi}^2}{g_{\phi\phi}}=\frac{r^2}{g_{\phi\phi}}.
\end{equation}
Consequently, the relative dust speed and differential-rotation parameter are
\begin{equation}\label{v_ell}
v=\frac{\sqrt{g_{\phi\phi}}(\Omega-\OmZ)}{N}=\frac{\eta}{r},\qquad
\ell(\eta)=\frac{H_{,\eta}}{H}.
\end{equation}
We work locally with $r>0$, $v\in C^2$, and $0<v<1$. For $\ell\equiv\ellzero\in\mathbb R$, one has
\begin{equation}\label{eq:constant-l-H}
H=-Ce^{\ellzero\eta},\qquad C>0.
\end{equation}
For $\ellzero\ne 0$ the angular velocity is
\begin{equation}\label{eq:constant-l-Omega}
\Omega(\eta)=\Omega_*-\frac{C\ellzero}{2}\operatorname{Ei}(\ellzero\eta),
\end{equation}
on an interval with $\eta>0$. For $\ellzero=0$, it is a constant $\Omega_0$. The constants are fixed by normalization and boundary data. With $\Drz v=v_{rr}+v_{zz}$ and $|\gradRZ v|^2=v_r^2+v_z^2$, the exact velocity equation is $\EE[v]=0$, where
\begin{align}\label{eq:exact-constant-l}
\EE[v]={}&\left[2+\ellzero r\left(\frac1v-v\right)\right]\Drz v
+\left[\frac2r+3\ellzero\left(\frac1v-v\right)\right]v_r\nonumber\\
&-\ellzero r\left(\frac1{v^2}+1\right)|\gradRZ v|^2
+\frac{2\ellzero}{r}-\frac{2v}{r^2}.
\end{align}
The retained-order equation used in the edge approximation of \cite{Astesiano2022EPJC} is $\ET[v]=0$, with
\begin{equation}\label{eq:truncated-operator}
\ET[v]=\left(2+\frac{\ellzero r}{v}\right)\Drz v
+\left(\frac2r+\frac{3\ellzero}{v}\right)v_r
-\frac{\ellzero r}{v^2}|\gradRZ v|^2+\frac{2\ellzero}{r}-\frac{2v}{r^2}.
\end{equation}
The exact difference is
\begin{equation}\label{eq:remainder-identity}
\EE[v]-\ET[v]=-\ellzero\left(rv\Drz v+3vv_r+r|\gradRZ v|^2\right).
\end{equation}
All results below concern the exact equation unless the retained-order equation is expressly specified. We begin by introducing an auxiliary potential that reduces the exact velocity equation to a linear elliptic equation. We then determine where this transformation admits a regular local inverse, allowing the velocity to be recovered from the potential. For a fixed scale $\eta_*>0$, we set
\begin{equation}\label{eq:constant-l-Phi}
\mathcal F=\Phi_{\ellzero}(r,\eta),\qquad
\Phi_{\ellzero}(r,\eta)=2\eta+\ellzero r^2\log\frac{\eta}{\eta_*}
-\frac{\ellzero}{2}\eta^2.
\end{equation}
The derivative of this map is the principal coefficient of the velocity equation
\begin{equation}\label{eq:Phi-eta-A}
A_{\ellzero}=\partial_\eta\Phi_{\ellzero}
=2+\frac{\ellzero r^2}{\eta}-\ellzero\eta
=2+\ellzero r\left(\frac1v-v\right).
\end{equation}
A straighforward application of the chain rule gives
\begin{equation}\label{eq:linearizing-intermediate}
\mathcal F_{rr}-r^{-1}\mathcal F_r+\mathcal F_{zz}
=A_{\ellzero}(\Drz\eta-r^{-1}\eta_r)
+\frac{4\ellzero r}{\eta}\eta_r-\ellzero\left(\frac{r^2}{\eta^2}+1\right)|\gradRZ\eta|^2.
\end{equation}
Substitution of $\eta=rv$ yields the exact identity
\begin{equation}\label{eq:linearizing-identity}
\mathcal L\mathcal F=r\EE[v],\qquad
\mathcal L=\partial_r^2-r^{-1}\partial_r+\partial_z^2.
\end{equation}
Thus, $\mathcal L\mathcal F=0$ is equivalent to the velocity equation after reconstruction on a sheet where $A_{\ellzero}\ne 0$. Moreover, the analytic inverse-function theorem applies on each such sheet. The operator is elliptic for $r>0$ and is formally self-adjoint for the measure $r^{-1}dr\,dz$. A change of $\eta_*$ adds a constant multiple of $r^2$ to $\mathcal F$, which leaves $\mathcal L\mathcal F$ unchanged.


\section{Density and the admissible inverse sheet}
\label{sec:constant-field-constraint}

The auxiliary equation provides a local construction wherever the transformation from the velocity to the potential admits a regular inverse. Physical admissibility further requires nonnegative dust density. We first derive the resulting density bounds and identify the boundary at which the inverse transformation develops a fold. We then establish analyticity on regular inverse sheets and examine the sensitivity of auxiliary elliptic continuation to perturbations of the prescribed data. Finally, we consider the additional restrictions imposed by extension across a complete regular axis under the stated global hypotheses. These results describe how density, regularity, and continuation constrain the local construction.


\subsection{Constant fields and density bounds}

A spatially constant speed already illustrates the distinction between
prescribing local data and solving the field equation on an open region.

\begin{proposition}\label{prop:constant-l-constraint}
Let $I\subset(0,\infty)$ and $J\subset\mathbb R$ be nonempty open
intervals. For constant $\ellzero$, no field $v(r,z)=v_0\neq0$ solves
the exact or retained-order constant-$\ell$ equation on $I\times J$.
\end{proposition}
\begin{proof}
All derivatives of $v_0$ are zero, so
\begin{equation}\label{eq:constant-residual}
\EE[v_0]=\ET[v_0]=\frac{2}{r^2}(\ellzero r-v_0).
\end{equation}
Vanishing on an interval requires $\ellzero r=v_0$ at two distinct
radii, which gives $\ellzero=0$ and then $v_0=0$.
\end{proof}
For $v_0/\ellzero>0$, the residual vanishes on the single cylinder
$r=v_0/\ellzero$. In particular, the calibration
$v_0=a\vc=\ellzero\RG$ used in~\cite{Astesiano2022EPJC} selects
$r=\RG$. Flat equatorial data leave the transverse curvature available
to satisfy the equation. Their local realization is established in
Theorem~\ref{thm:local-flat-trace-existence}. At a regular point with $r>0$, $0<v<1$, and $\eta=rv$, the exact density is
\begin{equation}\label{eq:exact-density}
8\pi G_N\rho=\frac{1}{4g_{rr}}
\left[\frac{\eta^2}{r^2}(2-\eta\ell)^2-r^2\ell^2\right]
\frac{\eta_r^2+\eta_z^2}{\eta^2}.
\end{equation}
Here, $g_{rr}=e^\mu>0$. Writing $x=r\ell$, we obtain
\begin{align}
\mathcal N_v(x)
&=4v^2-4v^3x-(1-v^4)x^2,\label{eq:density-polynomial}\\
\rho&=\frac{\eta_r^2+\eta_z^2}{32\pi G_N g_{rr}\eta^2}
\mathcal N_v(x).\label{eq:density-sign-decomposition}
\end{align}
The following proposition translates the requirement of nonnegative dust density into explicit bounds on the variables governing the inverse transformation.
\begin{proposition}\label{prop:density-positivity-domain}
At a regular point with $r>0$, $0<v<1$, and finite $\ell(\eta)$,
the algebraic factor $\mathcal N_v(r\ell)$ is nonnegative precisely when
\begin{equation}\label{eq:density-domain}
-\frac{2v}{1-v^2}\leqslant r\ell
\leqslant\frac{2v}{1+v^2}.
\end{equation}
If $\nabla_{rz}\eta\neq 0$, this condition is equivalent to
$\rho\geqslant 0$, and strict inequalities are equivalent to $\rho>0$.
Either endpoint gives $\rho=0$.
\end{proposition}

\begin{proof}
The factorization
\begin{equation}\label{eq:density-factorization}
\mathcal N_v(x)=[2v-(1+v^2)x][2v+(1-v^2)x]
\end{equation}
has the two roots
\begin{equation}\label{eq:density-roots}
x_-=-\frac{2v}{1-v^2},\qquad
x_+=\frac{2v}{1+v^2}.
\end{equation}
Its leading coefficient is negative, and $x_-<0<x_+$. Hence, it is nonnegative exactly on $[x_-,x_+]$ and positive in the interior. The remaining factor in \eqref{eq:density-sign-decomposition} is positive exactly when $\nabla_{rz}\eta\neq 0$.
\end{proof}
At a critical point of $\eta$, the density vanishes independently of
$\mathcal N_v$ while continuity controls this exception on an open
nonnegative-density region. Furthermore, the density formula also allows us to characterize the open regions in which the solution reduces to a vacuum plateau.
\begin{corollary}\label{cor:vacuum-plateau}
Let $U$ be an open regular region on which $0<v<1$, $\rho\geqslant 0$,
and the factors in \eqref{eq:density-sign-decomposition} are continuous.
If \eqref{eq:density-domain} fails at $p\in U$, then $p$ has a connected
open neighbourhood on which $\eta$ is constant and $\rho=0$. Consequently, the density interval holds throughout $U$ whenever $\eta$ has no open constant plateau.
\end{corollary}
\begin{proof}
Failure of the interval makes $\mathcal N_v(r\ell)<0$ near $p$. Nonnegative density then forces $\nabla_{rz}\eta=0$ throughout a connected neighbourhood. The density formula gives $\rho=0$ there.
\end{proof}
For an exactly flat trace the critical-point exception disappears.
\begin{corollary}\label{cor:flat-trace-density-radius}
Suppose $v(r,0)=v_0\in(0,1)$ on an interval and $\ell=\ellzero$ is constant. At regular trace points, nonnegative density is equivalent to
\begin{eqnarray}
0<r&&\leqslant R_\rho^{(+)}=\frac{2v_0}{\ellzero(1+v_0^2)}
\qquad\text{if }\ellzero>0,\label{eq:positive-l-cutoff}\\
0<r&&\leqslant R_\rho^{(-)}=\frac{2v_0}{|\ellzero|(1-v_0^2)}
\qquad\text{if }\ellzero<0.\label{eq:negative-l-cutoff}
\end{eqnarray}
The density vanishes at the stated radius and is negative beyond it.
For $\ellzero=0$, $\rho>0$ at every regular trace point. These conclusions require no condition on the transverse derivatives.
\end{corollary}
\begin{proof}
Along the trace, $\eta_r=v_0\neq 0$. Apply Proposition~\ref{prop:density-positivity-domain}, using the upper bound for $\ellzero>0$ and the lower bound for $\ellzero<0$. For $\ellzero=0$, the algebraic factor is $4v_0^2>0$.
\end{proof}
These radii characterize the sign of the density along the prescribed
trace. A material boundary must also satisfy the junction conditions
discussed below.

\subsection{The lower density boundary and analytic reconstruction}
\label{sec:analytic-global-structure}
For constant $\ellzero$, we set
\begin{equation}\label{A}
A=\partial_\eta\Phi_{\ellzero}(r,\eta)=2+\ellzero r\left(\frac1v-v\right).
\end{equation}
The density factor becomes
\begin{equation}\label{eq:density-fold-factorization}
\mathcal N_v(x)=vA[2v-(1+v^2)x],\qquad x=r\ellzero.
\end{equation}
Thus, $A=0$ precisely at the lower root $x_-$. This endpoint is a fold of the reconstruction map, where local invertibility fails. At the upper root, $A=4/(1+v^2)>0$, so the density vanishes while the reconstruction map retains a regular local inverse. The following theorem makes this distinction precise by showing that the loss of local invertibility at the lower density boundary has the structure of a Whitney fold.
\begin{theorem}\label{thm:density-whitney-fold}
At a point $(r_0,\eta_0)$ with $\eta_0=r_0v_0$, $0<v_0<1$, and
\begin{equation}\label{eq:lower-root-is-A-zero}
r_0\ellzero=-\frac{2v_0}{1-v_0^2},
\end{equation}
the map $\Psi_{\ellzero}(r,\eta)=(r,\Phi_{\ellzero}(r,\eta))$
has a nondegenerate Whitney fold. Strictly positive density selects
the inverse sheet $A>0$. An auxiliary datum transverse to the fold has
square-root inverse behaviour. A smooth inverse through the fold
requires the pulled-back discriminant to be the square of a smooth
signed function with a coherent sheet choice.
\end{theorem}
\begin{proof}
The Jacobian matrix and its determinant are
\begin{equation}\label{eq:point-linearising-jacobian}
D\Psi_{\ellzero}=\begin{pmatrix}1&0\\ \partial_r\Phi_{\ellzero}&A\end{pmatrix},
\qquad \det D\Psi_{\ellzero}=A.
\end{equation}
At the lower root, $\ellzero<0$ and
\begin{equation}\label{eq:fold-second-derivative}
\partial_{\eta\eta}\Phi_{\ellzero}
=-\ellzero\left(1+\frac1{v_0^2}\right)>0.
\end{equation}
The differential has rank one, its kernel is generated by
$\partial_\eta$, and the determinant has nonzero derivative in this
direction. The parameter-dependent Morse lemma therefore gives
coordinates in which $\Psi_{\ellzero}(R,S)=(R,S^2)$. This is the fold normal form~\cite{GolubitskyGuillemin1973Stable}. Positive density places $x$ strictly between $x_-$ and $x_+$, so both linear factors in \eqref{eq:density-factorization} are positive. Equation~\eqref{eq:density-fold-factorization} then yields $A>0$. Finally, inversion of the normal form gives $S=\pm\sqrt Y$. A transverse zero of the target coordinate $Y$ produces square-root behaviour. Moreover, a smooth inverse exists exactly when $Y$ has a compatible smooth signed square root in these coordinates.
\end{proof}
Along a one-dimensional slice, finite even order and a positive leading
coefficient ensure a smooth signed square root. In several variables,
the required condition is factorization as a square. However, after inversion, metric regularity must still be checked in the reconstruction formulas. The next result shows that on a regular inverse sheet, elliptic regularity yields analyticity, which in turn determines how a locally constant region extends through the connected domain.
\begin{theorem}\label{thm:automatic-analyticity}
Let $D\subset\{r>0\}$ be connected and let $\eta\in C^2(D)$ be positive.
Suppose $v=\eta/r$ solves the exact constant-$\ell$ equation and
$A\neq 0$ throughout $D$. Then, $\eta$, $v$, and all locally reconstructed
metric coefficients are real analytic. If $\eta$ is constant on a
nonempty open subset, it is constant throughout $D$.
\end{theorem}
\begin{proof}
The chain-rule identity \eqref{eq:linearizing-identity} gives
\begin{equation}
\mathcal L\mathcal F=0,\qquad
\mathcal F=\Phi_{\ellzero}(r,\eta),\qquad
\mathcal L=\partial_r^2-r^{-1}\partial_r+\partial_z^2.
\end{equation}
Uniform ellipticity on compact subdomains of $D$ and analyticity of the coefficients imply $\mathcal F\in C^\omega(D)$ by analytic elliptic regularity~\cite{MorreyNirenberg1957CPAM}. Since $\partial_\eta\Phi_{\ellzero}=A\neq 0$, the analytic implicit-function theorem transfers this regularity to $\eta$ and $v$. Analyticity then passes to the Killing block and to the closed one-form defining the meridional quadrature, whose local primitives are therefore analytic. Suppose now that $\eta=\eta_0>0$ on an open set and define
\begin{equation}\label{eq:constant-plateau-comparison}
\mathcal F_0(r)=\Phi_{\ellzero}(r,\eta_0)
=2\eta_0-\frac{\ellzero}{2}\eta_0^2
+\ellzero r^2\log\left(\frac{\eta_0}{\eta_*}\right).
\end{equation}
Since both $1$ and $r^2$ are $\mathcal L$-harmonic, the difference $\mathcal F-\mathcal F_0$ is $\mathcal L$-harmonic and vanishes on an open set. Unique continuation therefore gives $\mathcal F=\mathcal F_0$ throughout $D$~\cite{Aronszajn1957JMPA}. The identity $\mathcal F=\mathcal F_0$ and local uniqueness of the inverse at $A\neq 0$ imply that the nonempty set ${\eta=\eta_0}$ is open. This set is also closed by continuity, so connectedness of $D$ yields $\eta=\eta_0$ everywhere.
\end{proof}
Strictly positive density supplies the regular-sheet hypothesis through
$A>0$. The same argument explains why a vacuum plateau, once formed,
propagates across a connected regular sheet.

\subsection{Continuation and the elliptic Cauchy problem}

Automatic analyticity leaves a separate question of continuous
dependence on the prescribed equatorial data. The exact auxiliary
operator has an explicit sequence exhibiting failure of such
dependence. We next examine the stability of auxiliary elliptic continuation by constructing solutions whose Cauchy data become arbitrarily small while their amplitudes grow away from the initial surface.
\begin{theorem}\label{thm:bessel-hadamard}
For positive integers $k$, let
\begin{equation}\label{Fk}
\mathcal F_k(r,z)=e^{-\sqrt k}rJ_1(kr)\cosh(kz).
\end{equation}
These functions satisfy $\mathcal L\mathcal F_k=0$. On every compact
annulus $I=[r_1,r_2]\Subset(0,\infty)$, their Cauchy data at $z=0$
converge to zero in every fixed $C^m$ and Sobolev norm. For every
$z_0>0$, however,
\begin{equation}
\|\mathcal F_k(\cdot,z_0)\|_{L^2(I)}\longrightarrow\infty.
\end{equation}
Hence, continuation from smooth Cauchy data to a collar of positive
depth is discontinuous.
\end{theorem}
\begin{proof}
The Bessel equation implies
\begin{equation}
\left(\partial_r^2-r^{-1}\partial_r\right)[rJ_1(kr)]=-k^2rJ_1(kr),
\end{equation}
which cancels the second $z$ derivative of $\cosh(kz)$. At $z=0$ the normal datum vanishes. Every fixed derivative of the Dirichlet datum is bounded on $I$ by a polynomial in $k$ times $e^{-\sqrt k}$, giving convergence in the stated norms. Uniform large-argument asymptotics on $I$ give
\begin{equation}
rJ_1(kr)=\sqrt{\frac{2r}{\pi k}}\cos(kr-3\pi/4)+O(k^{-3/2})
\end{equation}
and therefore,
\begin{equation}
\|rJ_1(kr)\|_{L^2(I)}^2=\frac{1}{\pi k}\int_I r\,dr+O(k^{-2})
\end{equation}
as $k\to\infty$~\cite{NISTDLMF}. Multiplication by
$e^{-\sqrt k}\cosh(kz_0)$ proves divergence. A continuous solution
map would send the null sequence of data to a null sequence in the
interior, yielding a contradiction.
\end{proof}
A complementary approach is to approximate a local $\mathcal L$-harmonic field on a compact interior set by fields defined on a larger domain. Runge approximation permits this construction when the complement of the compact set in the larger domain has no relatively compact component. With sufficient control of derivatives, the approximation preserves strict margins for inverse-sheet regularity, subluminality, and density on a protected core. The precise statement and the hypotheses needed for reconstruction are given in the Supplementary Material. Since this control is confined to the protected core, the outer geometry remains subject to a separate completion problem.


\section{Local flat traces and regular-axis completion}\label{sec:exact-flat-traces}

We now examine whether a speed profile that is constant along an equatorial interval can arise from a solution with nontrivial transverse dependence. Using the auxiliary potential, we construct such traces on admissible annuli and reconstruct the corresponding analytic dust metrics with positive density. We then investigate whether these local realizations can extend to a regular rotation axis. Under the stated axis hypotheses, regularity forces $\ellzero=0$, after which analyticity propagates the constant equatorial speed towards the axis, contradicting its required decay. This leads us to consider continuation within annular domains, where an additional transverse-maximum condition yields an explicit finite radial window for an exactly flat trace.

\subsection{Analytic realization and metric reconstruction}

We begin by showing that an exactly flat equatorial trace can be realized by a local analytic dust spacetime whenever the prescribed data lie strictly within the admissible density interval.
\begin{theorem}\label{thm:local-flat-trace-existence}
Let $r_{*}>0$, $0<v_{0}<1$, and $\ellzero\in\mathbb R$, with
\begin{equation}\label{eq:noncharacteristic-flat-trace}
A_{*}=2+\ellzero r_{*}(1/v_{0}-v_{0})\neq 0.
\end{equation}
There is a neighbourhood
$D_{*}=(r_{*}-\delta,r_{*}+\delta)\times(-z_{*},z_{*})$, where
$0<\delta<r_{*}$, on which the exact constant-$\ell$ equation has a
unique real-analytic solution with Cauchy data
$v(r,0)=v_{0}$ and $v_z(r,0)=0$. The solution is even in $z$ and satisfies
\begin{equation}\label{eq:local-flat-trace-expansion}
v(r,z)=v_{0}+\frac{v_{0}-\ellzero r}{r^{2}[2+\ellzero r(1/v_{0}-v_{0})]}z^{2}+O(z^{4}).
\end{equation}
With $H=-Ce^{\ellzero\eta}$, the angular velocity $\Omega$ is obtained by integrating Eq.~\eqref{eq:model-velocity-relations}, while $\mu$ follows from the quadratures \eqref{eq:local-mu-quadrature-r}--\eqref{eq:local-mu-quadrature-z}. Fixing $C>0$, the additive constant in $\Omega$, and the reference value $\mu(r_*,0)$ therefore determines a unique real-analytic local metric of the exact rotating-dust family. If
\begin{equation}\label{eq:strict-density-at-rstar}
-\frac{2v_{0}}{1-v_{0}^{2}}<\ellzero r_{*}
<\frac{2v_{0}}{1+v_{0}^{2}},
\end{equation}
then $D_{*}$ can be chosen so that $0<v<1$ and $\rho>0$ throughout it.
\end{theorem}
\begin{proof}
Prescribe the analytic auxiliary data
\begin{equation}\label{eq:F-flat-boundary-data}\\
\mathcal F(r,0)=\Phi_{\ellzero}(r,rv_{0})=2rv_{0}+\ellzero r^{2}\log\!\left(\frac{rv_{0}}{\eta_{*}}\right)
-\frac{\ellzero}{2}r^{2}v_{0}^{2},\qquad
\mathcal F_z(r,0)=0.
\end{equation}
Since $r_{*}>0$, the equation $\mathcal F_{zz}=-\mathcal F_{rr}+r^{-1}\mathcal F_r$ has analytic coefficients near $(r_{*},0)$ and a noncharacteristic Cauchy surface $z=0$. The Cauchy--Kowalevski theorem gives a unique local analytic solution \cite{Evans2010PDE}. Reflection in $z$ preserves both the equation and its data, so uniqueness makes $\mathcal F$ even. The condition $A_{*}\neq 0$ is exactly
$\partial_\eta\Phi_{\ellzero}(r_{*},r_{*}v_{0})\neq 0$.
The analytic implicit-function theorem therefore recovers a unique
analytic $\eta$ with $\mathcal F=\Phi_{\ellzero}(r,\eta)$ while the inverse inherits reflection symmetry and the prescribed trace. Moreover, the linearizing identity proves that $v=\eta/r$ solves the exact equation. Let us now shrink the neighbourhood so that $\eta>0$ and $0<v<1$. The functions
$H=-Ce^{\ellzero\eta}$ and $\Omega$, with
$\Omega_{,\eta}=H_{,\eta}/(2\eta)$, are analytic there.
Equation~\eqref{eq:model-metric-components} defines an analytic Killing
block $G$ satisfying $g_{\phi\phi}>0$ and $\det G=-r^{2}$.
To reconstruct the remaining metric function, we write
$A=\partial_\eta\Phi_{\ellzero}=2+\ellzero r^{2}/\eta-\ellzero\eta$.
The exact quadratures \cite{Winicour1975JMP,Astesiano2022EPJC} reduce to
\begin{align}
M_r=\mu_r
&=\frac{(g_{tt})_r(g_{\phi\phi})_r-(g_{tt})_z(g_{\phi\phi})_z
-(g_{t\phi})_r^2+(g_{t\phi})_z^2}{2r}=\frac{A^2}{8r}(\eta_z^2-\eta_r^2),
\label{eq:local-mu-quadrature-r}\\
M_z=\mu_z&=\frac{(g_{tt})_z(g_{\phi\phi})_r+(g_{tt})_r(g_{\phi\phi})_z
-2(g_{t\phi})_r(g_{t\phi})_z}{2r}
=-\frac{A^2}{4r}\eta_r\eta_z.
\label{eq:local-mu-quadrature-z}
\end{align}
For completeness, the simplification can be checked at each point after
a constant change of angular coordinate sets $\Omega$ to zero there.
Then, $d g_{tt}=(H_{,\eta}-2\eta\Omega_{,\eta})d\eta=0$ and
$d g_{t\phi}=(A/2)d\eta$. Both quadrature expressions are invariant
under this constant change of Killing basis. Their compatibility follows
from the explicit identity
\begin{equation}\label{eq:local-mu-integrability}
\partial_zM_r-\partial_rM_z
=\frac{A\eta_z}{4r}\mathcal L\mathcal F=0.
\end{equation}
Thus, $M_r\,dr+M_z\,dz$ is a closed analytic one-form on a sufficiently
small rectangle. Integrating it with the assigned value at $(r_{*},0)$
gives a unique analytic $\mu$. The resulting line element
$G_{AB}dx^A dx^B+e^\mu(dr^2+dz^2)$ is Lorentzian and satisfies the
remaining equations of the exact dust construction. Along the trace, we have $v_r=v_{rr}=v_z=0$. The field equation gives
\begin{equation}\label{eq:exact-flat-trace-balance}
[2+\ellzero r(1/v_{0}-v_{0})]v_{zz}(r,0)
+\frac{2\ellzero}{r}-\frac{2v_{0}}{r^{2}}=0,
\end{equation}
and hence
\begin{equation}\label{eq:exact-flat-trace-curvature}
v_{zz}(r,0)=
\frac{2(v_{0}-\ellzero r)}{r^{2}[2+\ellzero r(1/v_{0}-v_{0})]}.
\end{equation}
Evenness gives the stated expansion. To establish positivity of the density, observe that the prescribed data yield
$\eta_r(r_{*},0)=v_{0}>0$ and $\eta_z(r_{*},0)=0$. The strict bounds \eqref{eq:strict-density-at-rstar} therefore imply $\rho(r_{*},0)>0$ by Proposition~\ref{prop:density-positivity-domain}. Finally, continuity preserves this inequality and subluminality on a smaller patch.
\end{proof}
This construction fixes the complete local geometry. Extending the
resulting annular patch requires additional information about the axis,
the matter boundary, and the exterior region. We first examine the axis.

\subsection{The metric condition at a regular axis}

Let $k$ and $m$ be the stationary and axial Killing fields. A regular
axis is a fixed component of the effective $U(1)$ action generated by
$m$, with angular period $2\pi$. In transverse Cartesian coordinates, $m=-y\partial_x+x\partial_y$, and $\mathcal A=\{x=y=0\}$. The next result combines metric regularity at the axis with the exact dust identities to establish $H'(0)=0$, which forces $\ellzero=0$ in the constant-$\ell$ class.

\begin{theorem}\label{thm:metric-axis-jet-constraint}
Let a stationary Lorentz metric extend as a $C^2$ metric across a regular
axis component $\mathcal A$. Assume that $k$ commutes with $m$ and is
timelike on $\mathcal A$. Let a nontrivial $(\eta,H)$ branch meet the
axis and satisfy the exact relations
\eqref{eq:model-velocity-relations}--\eqref{eq:model-metric-components}
away from it, with $0<|\eta|<r$ and $\eta\to 0$ at the axis. If $H$ is $C^1$ at zero and $H(0)=-C<0$, then $H'(0)=0$. In particular, a constant-$\ell$ branch satisfying these hypotheses has
$\ellzero=0$.
\end{theorem}

\begin{proof}
Let $\varrho$ be the transverse distance to the axis. The $C^2$ metric
and rotational invariance give the uniform local expansions
\begin{equation}\label{eq:metric-axis-jets}
g(m,m)=a\varrho^2+o(\varrho^2),\qquad
g(k,m)=O(\varrho^2),\qquad
g(k,k)=-D+O(\varrho^2),
\end{equation}
where $a,D>0$ may depend on the coordinate along the axis. The determinant identity yields
\begin{equation}\label{eq:axis-r-rho-comparison}
r^2=Da\varrho^2+o(\varrho^2),\qquad r\asymp\varrho.
\end{equation}
Subluminality therefore gives $\eta\to 0$ along the branch. The differentiated axis expansions are
\begin{equation}\label{eq:metric-axis-derivative-jets}
\partial_\varrho g(k,k)=O(\varrho),\qquad
\partial_\varrho g(k,m)=O(\varrho),\qquad
\partial_\varrho g(m,m)=2a\varrho+o(\varrho).
\end{equation}
Here, the derivative fixes the coordinate along the axis. These estimates
hold uniformly near any fixed axis point. If we suppose $h_1=H'(0)\neq0$,  integration of $\Omega_{,\eta}=H_{,\eta}/(2\eta)$ gives
\begin{equation}
\Omega(\eta)=\frac{h_1}{2}\log|\eta|+o\bigl(|\log|\eta||\bigr),
\end{equation}
so $|\Omega|\to\infty$ on approaching the axis through the branch.
On the other hand, differentiating
$H=g(k,k)+2\Omega g(k,m)+\Omega^2g(m,m)$ and using
$g(k,m)+\Omega g(m,m)=\eta$ gives
\begin{equation}\label{eq:axis-geodesic-differential-identity}
dg(k,k)+2\Omega\,dg(k,m)+\Omega^2\,dg(m,m)=dH-2\eta\,d\Omega=0.
\end{equation}
If we take the transverse radial component and divide by
$\varrho\Omega^2$. the first two terms tend to zero by
\eqref{eq:metric-axis-derivative-jets}, whereas the third tends to
$2a>0$. This contradiction proves $H'(0)=0$. For $H=-Ce^{\ellzero\eta}$, the conclusion is $C\ellzero=0$.
\end{proof}
The quadratic axis expansion is essential here because elementary flatness on a punctured neighbourhood alone does not supply it. After a time
normalization with $H(0)=-1$, an analytic $H$ must consequently satisfy
$H(\eta)=-1+O(\eta^2)$. The same coefficient condition was obtained
under additional material assumptions in \cite{Galoppo2025EPJC}.
A complementary proof based on the regularity of the dust congruence is
given in the supplementary material. Having established that regular-axis completion requires $\ellzero=0$, we now determine the rate at which the ZAMO speed must vanish as the axis is approached.
\begin{lemma}\label{lem:rigid-axis-v-zero}
Under the metric and axis hypotheses of
Theorem~\ref{thm:metric-axis-jet-constraint}, a branch with
$\ellzero=0$ satisfies $\eta=O(r^2)$, and $v=O(r)$ at the axis.
\end{lemma}
\begin{proof}
The angular velocity is a finite constant $\Omega_0$, and $\eta=g(k,m)+\Omega_0g(m,m)$. Both terms on the right are $O(\varrho^2)$ by
\eqref{eq:metric-axis-jets}. Since $r\asymp\varrho$, division by $r$
gives the asserted decay.
\end{proof}

\subsection{A complete regular-axis branch}

To study global rigidity on the $\ellzero=0$ branch, we recast the auxiliary equation $\mathcal L\mathcal F=0$ as a harmonic problem. The following proposition shows that the substitution $\mathcal F=r^2U$ yields the five-dimensional Laplace equation away from the axis, providing the basis for the global argument.
\begin{proposition}\label{prop:five-dimensional-lift}
We write $\mathcal F=r^2U$ for $r>0$ and set
$\widetilde U(y,z)=U(|y|,z)$ for $y\in\mathbb R^4$. Then
\begin{equation}\label{eq:five-dimensional-lift}
\mathcal L\mathcal F=r^2\left(U_{rr}+\frac3rU_r+U_{zz}\right)
=r^2\Delta_{\mathbb R^5}\widetilde U.
\end{equation}
Thus, $\mathcal L$-harmonic fields correspond to harmonic fields in
five dimensions invariant under rotations of $y$.
\end{proposition}
\begin{proof}
Differentiating $r^2U$ gives $\partial_r(r^2U)=2rU+r^2U_r$ and
$\partial_r^2(r^2U)=2U+4rU_r+r^2U_{rr}$. Moreover, substitution yields the first identity. The radial Laplacian in $\mathbb R^4_y$ is $\partial_r^2+3r^{-1}\partial_r$, which gives the second.
\end{proof}
The harmonic lift allows us to combine regularity along the complete axis with decay at infinity to show that, under the hypotheses below, the $\ellzero=0$ branch reduces to Minkowski spacetime.
\begin{theorem}\label{thm:ellzero-liouville}
Suppose the $\ellzero=0$ branch is defined for all $r>0$ and
$z\in\mathbb R$ and extends regularly to the complete axis.
Assume there is no finite boundary, puncture, additional end,
horizon, or singular axis segment. For every compact interval
$K\Subset\mathbb R$, assume
\begin{equation}
|\eta(r,z)|\leqslant C_Kr^2
\quad\text{for }z\in K,\quad 0<r<r_K
\end{equation}
with suitable positive constants $C_K$ and $r_K$. Suppose also that
\begin{equation}\label{eq:eta-outer-decay}
|\eta(r,z)|\leqslant C_\infty\frac rR,
\qquad R=\sqrt{r^2+z^2}\geqslant R_0.
\end{equation}
Then, $\eta=0$. After stationary and meridional normalization,
the metric is Minkowski and the dust density vanishes.
\end{theorem}
\begin{proof}
For $\ellzero=0$, one has $\mathcal F=2\eta$, so
\begin{equation}\label{eq:lifted-eta-field}
\widetilde U(y,z)=\frac{2\eta(|y|,z)}{|y|^2}
\end{equation}
is harmonic away from the line
$\Gamma=\{(0,z)\in\mathbb R^4\times\mathbb R\}$.
The axis estimate bounds $\widetilde U$ near every compact segment
of $\Gamma$. Such segments have zero Newtonian capacity in
$\mathbb R^5$, so the removable-singularity theorem extends
$\widetilde U$ harmonically across $\Gamma$. On a sphere of radius $S$, let $r=S\sin\theta$. The outer estimate gives
\begin{equation}\label{eq:liouville-pointwise-bound}
|\widetilde U|\leqslant\frac{2C_\infty}{S^2\sin\theta}.
\end{equation}
The angular measure on $\mathbb S^4$ contains
$\sin^3\theta\,d\theta$, so this bound is integrable at both poles.
Consequently, we find
\begin{equation}\label{eq:liouville-mean-decay}
\frac1{|\mathbb S^4_S|}\int_{\mathbb S^4_S}|\widetilde U|\,d\sigma
\leqslant\frac C{S^2}.
\end{equation}
For fixed $X\in\mathbb R^5$ and $S>2|X|$, the Poisson kernel on
the ball of radius $S$ is bounded by a fixed multiple of the
normalized spherical measure. Furthermore, the Poisson formula and
\eqref{eq:liouville-mean-decay} imply
$|\widetilde U(X)|\leqslant C_XS^{-2}$.
Letting $S\to\infty$ gives $\widetilde U=0$, and hence $\eta=0$.
The remaining relations give constant $H=-C$, constant $\Omega$,
$g_{\phi\phi}=r^2/C$, and constant $\mu$. Regularity of the axis
and the stated normalizations yield Minkowski spacetime.
\end{proof}
The theorem concerns the complete orbit half-plane with the stated
axis and decay estimates. A finite body matched to a vacuum exterior
requires separate analysis. Notice that within the strictly positive-speed class, the theorem excludes a branch satisfying all these global hypotheses. We now combine analyticity with the required axis decay to show that a positive exactly flat equatorial trace cannot extend along a connected regular branch to the rotation axis.
\begin{theorem}\label{thm:all-constant-l-flat-trace-no-go}
Let an exact constant-$\ell$ solution satisfy the hypotheses of
Theorem~\ref{thm:metric-axis-jet-constraint}. Assume that its equatorial
radial domain
$\mathcal I_0=\{r>0\mid(r,0)\in\mathcal U\}$ is a connected open
interval whose closure meets the regular axis at $r=0$.
Then, no nonempty open interval $I\subset\mathcal I_0$ supports $v(r,0)=v_0$ with $0<v_0<1$.
\end{theorem}
\begin{proof}
The metric axis condition forces $\ellzero=0$, so
$\mathcal F=2\eta$ and $\eta_{rr}-r^{-1}\eta_r+\eta_{zz}=0$. This equation is elliptic with analytic coefficients on $r>0$. 
Interior analytic regularity makes $\eta$, and therefore $v=\eta/r$,
real analytic there \cite{MorreyNirenberg1957CPAM}. The identity theorem applied to $v(r,0)-v_0$ propagates the constant trace across the connected interval $\mathcal I_0$. It follows that $v(r,0)\to v_0>0$ as $r\downarrow 0$. Lemma~\ref{lem:rigid-axis-v-zero} instead gives $v(r,0)=O(r)$,
which proves the claim.
\end{proof}
Connectedness ensures that the flat interval and the axis lie on one
analytic equatorial branch. Even though a junction, a vacuum gap, or a singularity can spoil this argument, the theorem requires no reflection
symmetry, density sign, or vertical-curvature assumption. For $\ellzero=0$, the reflection-symmetric local continuation is
$v=v_0\sqrt{r^2+z^2}/r$. Its direct verification and finite luminal
height are recorded in the supplementary material. The next result shows that, for an annular trace separated from the axis, nonnegative density and an equatorial maximum in the transverse direction together restrict the radial interval over which an exactly flat speed profile can persist.
\begin{theorem}\label{thm:annular-flat-trace-bound}
Let an exact constant-$\ell$ solution be $C^2$ on
$I\times(-z_*,z_*)$, where $I\subset(0,\infty)$ is a nonempty open
interval. Suppose $v(r,0)=v_0\in(0,1)$ and $\rho(r,0)\geq0$ on $I$.
Assume that $z=0$ is a local vertical maximum of $v$ at every $r\in I$.
Then $\ellzero>0$ and
\begin{equation}\label{eq:annular-radius-window}
\frac{v_0}{\ellzero}\leq r\leq
\frac{2v_0}{\ellzero(1+v_0^2)}.
\end{equation}
Moreover, strictly positive density and $v_{zz}(r,0)<0$ give strict inequalities.
\end{theorem}
\begin{proof}
Flatness and the maximum condition give $\eta_r=v_0>0$, $\eta_z=0$,
and $v_{zz}\leq 0$ on the trace. Every trace point is therefore
noncritical. Let us set $x=\ellzero r$. The density criterion yields
\begin{equation}\label{eq:flat-trace-x-domain}
x_-=-\frac{2v_0}{1-v_0^2}\leq x\leq
\frac{2v_0}{1+v_0^2}=x_+.
\end{equation}
The coefficient in the trace equation is
\begin{equation}\label{eq:A-of-x-flat-trace}
A(x)=2+\frac{1-v_0^2}{v_0}x.
\end{equation}
At $x=x_-$ it vanishes, whereas
\eqref{eq:exact-flat-trace-balance} would then require $x_-=v_0$.
This is impossible. Hence, $x>x_-$ and $A(x)>0$, so
\begin{equation}\label{eq:vzz-x-form}
v_{zz}(r,0)=\frac{2(v_0-x)}{r^2A(x)}\leq 0.
\end{equation}
It follows that $x\geq v_0>0$ and $\ellzero>0$. Combining this with $x\leq x_+$ proves the window. Positive density excludes $x=x_+$, while strictly negative vertical curvature gives $x>v_0$.
\end{proof}
For any $r_1<r_2$ in the trace interval
\begin{equation}\label{eq:flat-trace-span-ratio}
\frac{r_2}{r_1}\leq\frac{2}{1+v_0^2}<2.
\end{equation}
The bound uses only equatorial data and the sign of the transverse
curvature. In the next section, near-flat traces lead to quantitative
bounds on the same radial extent.


\section{Finite radial spans of equatorial traces}
\label{sec:approx-flat-constraint}

We now extend the analysis of exactly flat traces to profiles with a bounded logarithmic slope. For reflection-symmetric solutions, we express transverse concavity invariantly through the relative dust--ZAMO rapidity and derive an exact equatorial identity linking this condition to the evolution of the slope. Combining that identity with nonnegative density yields finite radial-span bounds whose form depends on the sign of $\ellzero$. We then construct analytic solutions establishing sharpness for the nonpositive class and for the positive class in the exactly flat limit. For positive $\ellzero$ and nonzero slope tolerance, the estimates provide upper bounds whose optimal value remains open. In every case, the bounds exclude a semi-infinite approximately flat trace under the stated hypotheses.


\subsection{Transverse shape and the logarithmic identity}

Let $\mathsf q=e^\mu(dr^2+dz^2)$ denote the orbit metric and let $r=\sqrt{-\det G}$ be the orbit-area radius. Where $dr\neq 0$, we define the unit vector field
\begin{equation}
N=\frac{J\operatorname{grad}_{\mathsf q}r}
{\lvert\operatorname{grad}_{\mathsf q}r\rvert_{\mathsf q}},
\qquad \psi=\operatorname{artanh}v,
\end{equation}
where $J$ is the orientation-compatible quarter turn. On the equatorial reflection surface, $N$ is transverse to that surface. We also impose the weak vertical-shape condition $N(N\psi)\leqslant 0$. While this constraint concerns the rapidity relative to ZAMOs, it imposes no independent requirement on the transverse profile of the density. Reflection symmetry allows us to express the invariant rapidity condition directly in terms of the transverse second derivative of the ZAMO speed, as the following proposition shows.
\begin{proposition}\label{prop:invariant-vertical-shape}
For a reflection-symmetric field with $0<v<1$, the identity
\begin{equation}\label{eq:invariant-vertical-shape}
N(N\psi)=e^{-\mu}\frac{v_{zz}}{1-v^2}
\end{equation}
holds on the equatorial plane. Thus, the vertical-shape condition is equivalent to $v_{zz}(r,0)\leqslant 0$.
\end{proposition}
\begin{proof}
In the Weyl chart, $\operatorname{grad}_{\mathsf q}r=e^{-\mu}\partial_r$ and its norm is $e^{-\mu/2}$. Hence, $N=e^{-\mu/2}\partial_z$, with the chosen orientation, and
\begin{equation}
N(N\psi)=e^{-\mu}\psi_{zz}+e^{-\mu/2}\partial_z(e^{-\mu/2})\psi_z.
\end{equation}
Reflection symmetry gives $v_z=\psi_z=0$. Using
$\psi_{zz}=v_{zz}/(1-v^2)+2v v_z^2/(1-v^2)^2$ and $v_z=0$ on the equatorial plane gives the stated formula. Since $e^{-\mu}/(1-v^2)>0$, the condition $N(N\psi)\leqslant 0$ is therefore equivalent to $v_{zz}(r,0)\leqslant 0$.
\end{proof}
Strict negativity gives a strict local transverse maximum. At a point where the second derivative vanishes, the weak condition leaves higher transverse derivatives unrestricted. This distinction matters because the solutions attaining the nonpositive-$\ellzero$ span bound have zero transverse curvature. To quantify departures from an exactly flat trace, we bound the logarithmic slope of the equatorial speed.
\begin{definition}\label{def:epsilon-flat}
Let $u(r)=v(r,0)>0$ on $[R_1,R_2]$, where $0<R_1<R_2$. For $0\leqslant\varepsilon<1$, the trace is $\varepsilon$-flat if
\begin{equation}\label{eq:epsilon-flat-definition}
\left|\frac{d\log u}{d\log r}\right|\leqslant\varepsilon.
\end{equation}
Its logarithmic radial span is
\begin{equation}\label{eq:logarithmic-span}
\Lambda=\log(R_2/R_1).
\end{equation}
\end{definition}
Notice that this definition is unchanged by a constant rescaling of the radial coordinate. For $R_1\leqslant r_1<r_2\leqslant R_2$, integration gives
\begin{equation}\label{eq:epsilon-flat-multiplicative-control}
(r_2/r_1)^{-\varepsilon}\leqslant
\frac{u(r_2)}{u(r_1)}\leqslant(r_2/r_1)^\varepsilon.
\end{equation}
In the following, we use the logarithmic variable $\sigma=\log(r/R_1)$ and write
\begin{equation}\label{eq:approx-flat-variables}
s=\frac{d\log u}{d\sigma},\qquad
\dot s=\frac{ds}{d\sigma},\qquad
y=\frac{\ellzero r}{u},\qquad
Z=\frac{r^2v_{zz}(r,0)}{u},\qquad
A=2+y(1-u^2).
\end{equation}
Then $|s|\leqslant\varepsilon$, and direct differentiation yields
\begin{equation}\label{eq:approx-flat-derivative-identities}
u_r=\frac{us}{r},\qquad
u_{rr}=\frac{u}{r^2}(\dot s+s^2-s),\qquad
\dot y=(1-s)y.
\end{equation}
In particular, $\eta_r(r,0)=u(1+s)>0$. Every point of such a trace is therefore noncritical for $\eta$, so the density interval applies everywhere along it. Expressing the exact velocity equation in these logarithmic variables relates the evolution of the equatorial slope to the transverse curvature and leads to a differential inequality under the density and concavity assumptions.
\begin{lemma}\label{lem:exact-equatorial-log-identity}
Let $v$ be a $C^2$ reflection-symmetric solution of the exact constant-$\ell$ VFE near $[R_1,R_2]\times\{0\}$, with $0<u<1$. Its equatorial data satisfy
\begin{equation}\label{eq:exact-equatorial-log-identity}
A(Z+\dot s)+2(1-yu^2)s^2+2y(1-u^2)s+2y-2=0.
\end{equation}
If the trace is $\varepsilon$-flat and $\rho(r,0)\geqslant 0$, then
\begin{equation}\label{eq:y-density-domain}
-\frac{2}{1-u^2}\leqslant y\leqslant\frac{2}{1+u^2},\qquad A>0.
\end{equation}
Under the vertical-shape condition, it follows that
\begin{equation}\label{eq:master-approx-flat-inequality}
A\dot s\geqslant
2(1-y)-2(1-yu^2)s^2-2y(1-u^2)s.
\end{equation}
\end{lemma}
\begin{proof}
Reflection symmetry gives $v_z(r,0)=0$. Substituting the derivative relations \eqref{eq:approx-flat-derivative-identities} into $\EE[v]=0$, with $\EE$ defined in \eqref{eq:exact-constant-l}, and multiplying by $r^2/u$ yields \eqref{eq:exact-equatorial-log-identity}. Since $\eta_r=u(1+s)>0$, Proposition~\ref{prop:density-positivity-domain} gives the stated interval for $y$. That interval implies $A\geqslant 0$. Equality could occur only for $y=-2/(1-u^2)$. At such a point, the logarithmic identity reduces to
\begin{equation}\label{eq:fold-boundary-substitutions}
(s+1)\bigl[(1+u^2)s-(3-u^2)\bigr]=0.
\end{equation}
The two possible slopes are $-1$ and $(3-u^2)/(1+u^2)>1$. Both contradict $|s|\leqslant\varepsilon<1$, so $A>0$. Finally, the vertical-shape condition gives $Z\leqslant 0$. Rearranging the identity yields
\begin{equation}\label{eq:equatorial-identity-rearranged}
A\dot s=2(1-y)-2(1-yu^2)s^2-2y(1-u^2)s-AZ,
\end{equation}
where the last term is nonnegative.
\end{proof}
The nonnegative-density and transverse-concavity hypotheses of the preceding lemma play distinct roles in the reduction. Density confines $y$ to a bounded interval while transverse concavity controls the evolution of the logarithmic slope. Their combination produces bounds independent of the inner-edge strength of the branch. A density-only estimate, with an explicitly calibrated inner-edge parameter, is given in the supplementary material.

\subsection{Universal bounds and their sharp cases}
Integrating the differential inequality from the preceding lemma yields explicit bounds on the logarithmic radial span, with separate estimates for nonpositive and positive $\ellzero$.
\begin{theorem}\label{thm:quantitative-approx-flat-no-go}
Let $v$ be a $C^2$ reflection-symmetric exact constant-$\ell$ solution near $[R_1,R_2]\times\{0\}$. Suppose its equatorial trace satisfies $0<u<1$, $\rho(r,0)\geqslant 0$, $N(N\psi)\leqslant 0$, and $|d\log u/d\log r|\leqslant\varepsilon$ for some $0\leqslant\varepsilon<1$. Set
\begin{equation}\label{eq:vmin-approx-flat}
\underline v=\min_{[R_1,R_2]}u>0.
\end{equation}
If $\ellzero\leqslant 0$, then
\begin{equation}\label{eq:approx-flat-bound-nonpositive-l}
\Lambda\leqslant2\operatorname{artanh}\varepsilon
=\log\frac{1+\varepsilon}{1-\varepsilon}.
\end{equation}
If $\ellzero>0$, then every $\alpha\in(0,1-\varepsilon)$ gives $\Lambda\leqslant\mathcal B_+(\varepsilon,\alpha,\underline v)$ with 
\begin{equation}\label{eq:approx-flat-bound-positive-l}
\mathcal B_+(\varepsilon,\alpha,\underline v)=\frac{\varepsilon(2+\alpha)}
{(1+\varepsilon)(1-\varepsilon-\alpha)}+\frac{1}{1-\varepsilon}
\log\frac{2}{(1+\underline v^2)\alpha}.
\end{equation}
\end{theorem}
\begin{proof}
The previous lemma gives $A>0$ and the differential inequality
\eqref{eq:master-approx-flat-inequality}. For $y\leqslant 0$, its right-hand side has the exact factorization
\begin{equation}\label{eq:sharp-negative-l-factorization}
2(1-y)-2(1-yu^2)s^2-2y(1-u^2)s-A(1-s^2)=-y(1+s)\bigl[3-u^2-(1+u^2)s\bigr].
\end{equation}
All factors on the right are nonnegative. The last factor is positive because $u<1$ and $s<1$. Hence, $\dot s\geqslant1-s^2$. Since $|s|<1$, it follows that
\begin{equation}
\frac{d}{d\sigma}\operatorname{artanh}s=\frac{\dot s}{1-s^2}\geqslant1.
\end{equation}
Integrating over $[0,\Lambda]$ gives
\begin{equation}
\Lambda\leqslant\operatorname{artanh}s(\Lambda)-\operatorname{artanh}s(0)
\leqslant 2\operatorname{artanh}\varepsilon.
\end{equation}
This proves the first bound, including the case $\varepsilon=0$, when a nondegenerate interval is excluded. For $\ellzero>0$, we have $y>0$ and
\begin{equation}\label{eq:y-growth-approx-flat}
\frac{d\log y}{d\sigma}=1-s\geqslant1-\varepsilon>0.
\end{equation}
Thus, $y$ increases strictly, while the density interval gives
\begin{equation}\label{eq:y-upper-vmin}
y\leqslant Y_{\max}=\frac{2}{1+\underline v^2}<2.
\end{equation}
Fix $\alpha\in(0,1-\varepsilon)$. On the initial portion where $y\leqslant\alpha$, we have $0<1-yu^2\leqslant 1$ and $A\leqslant 2+\alpha$. Consequently, we find that
\begin{equation}
A\dot s\geqslant 2(1-\alpha)-2\varepsilon^2-2\alpha\varepsilon
=2(1+\varepsilon)(1-\varepsilon-\alpha)>0,
\end{equation}
and hence, we have
\begin{equation}\label{eq:q-lower-small-y}
\dot s\geqslant c_\alpha:=
=\frac{2(1+\varepsilon)(1-\varepsilon-\alpha)}{2+\alpha}.
\end{equation}
Since $s$ remains in an interval of width $2\varepsilon$, this initial portion has logarithmic length at most $2\varepsilon/c_\alpha$. On the remaining portion, $y$ starts at a value at least $\alpha$ and grows logarithmically at rate at least $1-\varepsilon$. Its upper bound therefore limits the length of that portion to $(1-\varepsilon)^{-1}\log{(Y_{\max}/\alpha)}$. The same argument applies if either portion is empty. Adding the two lengths proves Eq.~\eqref{eq:approx-flat-bound-positive-l}.
\end{proof}
Both estimates are finite for every fixed $\varepsilon<1$. The nonpositive branch permits progressively shorter intervals as the slope tolerance decreases. Moreover, the positive branch retains a nonzero exact-flat limit, which agrees with the annular window obtained from the constant trace equation. The next result shows that an explicit analytic solution with $\ellzero=0$ attains the bound for the nonpositive class, while every solution with $\ellzero<0$ satisfies a strict inequality on each nondegenerate interval.
\begin{corollary}\label{cor:sharp-nonpositive-l-span}
For every $0<\varepsilon<1$, equality in
Eq.~\eqref{eq:approx-flat-bound-nonpositive-l} is attained by a real-analytic solution with $\ellzero=0$, positive density, and $v_{zz}=0$. For $\ellzero<0$, the inequality is strict on every nondegenerate interval.
\end{corollary}
\begin{proof}
Let us set $r=r_0 e^t$ and choose $u(t)=a\cosh t$, $s(t)=\tanh t$, and $|t|\leqslant\operatorname{artanh}\varepsilon$. The field
\begin{equation}
\eta(r,z)=ru(r)=\frac a2\left(\frac{r^2}{r_0}+r_0\right)
\end{equation}
is independent of $z$ and satisfies $\mathcal L(2\eta)=0$. Since $\mathcal F=2\eta$ when $\ellzero=0$, it defines an exact solution with zero transverse curvature. Its gradient is nonzero, so the density is positive. Choosing $0<a<\sqrt{1-\varepsilon^2}$ ensures $u<1$ throughout the interval and in a neighbourhood of it. Moreover, the logarithmic slope increases from $-\varepsilon$ to $\varepsilon$ over an interval of logarithmic length $2\operatorname{artanh}\varepsilon$, so the radial-span bound is attained. For $\ellzero<0$, the right-hand side of Eq.~\eqref{eq:sharp-negative-l-factorization} is strictly positive. Thus, $\dot s>1-s^2$, and integration gives a strict span inequality.
\end{proof}
The supplementary material also constructs strictly negative branches approaching the same supremum as $\ellzero\uparrow 0$. This approximation takes place within analytic solutions of the full two-dimensional field equation. For $\ellzero>0$, the parameter $\alpha$ can be optimized. An explicit minimizer and the small-$\varepsilon$ expansion are given in the supplementary material. The exact-flat limit already follows by taking, for example, $\alpha=1-\sqrt{\varepsilon}$ for sufficiently small positive $\varepsilon$. In this choice, the first term in $\mathcal B_+$ tends to zero, while the second tends to $\log[2/(1+\underline v^2)]$. For $\varepsilon=0$, taking $\alpha\uparrow 1$ directly in the theorem gives
\begin{equation}\label{eq:optimized-exact-flat-limit}
\frac{R_2}{R_1}\leqslant\frac{2}{1+v_0^2}
\end{equation}
when $u=v_0$. The next proposition shows that this limiting value is the supremum of the radial spans attained by local analytic solutions.
\begin{proposition}\label{prop:sharp-positive-l-exact-flat-limit}
Fix $0<v_0<1$ and $\ellzero>0$. Among reflection-symmetric real-analytic local solutions with $u=v_0$, strictly positive density, and $v_{zz}(r,0)<0$, the radial span has supremum
\begin{equation}\label{eq:positive-l-exact-flat-supremum}
\sup\frac{R_2}{R_1}=\frac{2}{1+v_0^2}.
\end{equation}
Strict positivity and strict transverse concavity exclude attainment on a closed interval.
\end{proposition}
\begin{proof}
For $s=\dot s=0$, Eq.~\eqref{eq:exact-equatorial-log-identity} gives $AZ=2(1-y)$. Thus, $v_{zz}<0$ is equivalent to $y>1$. The strict density interval gives $y<2/(1+v_0^2)$. Every admissible interval therefore lies inside
\begin{equation}
\frac{v_0}{\ellzero}<r<\frac{2v_0}{\ellzero(1+v_0^2)},
\end{equation}
which proves the upper bound and excludes equality. Choose a sequence of closed intervals contained in this open window, with their left and right endpoints approaching the corresponding boundaries of the window. On each interval, we prescribe the analytic auxiliary Cauchy data
\begin{equation}
\mathcal F(r,0)=\Phi_{\ellzero}(r,rv_0),\qquad \mathcal F_z(r,0)=0.
\end{equation}
The analytic construction of Theorem~\ref{thm:local-flat-trace-existence} applies at every point. Local uniqueness identifies the solutions on overlaps, and compactness gives a neighbourhood of the entire interval. Since $A>0$, the same inverse sheet recovers a reflection-symmetric analytic field. The strict density and transverse inequalities hold along its trace, while subluminality and positive density persist on a smaller neighbourhood. Finally, the resulting endpoint ratios tend to $2/(1+v_0^2)$.
\end{proof}
For positive $\ellzero$ and fixed $\varepsilon>0$, the universal comparison bound is not sharp. The supplementary material derives a strictly smaller bound by minimizing the exact slope drift over a state enclosure. The exact supremum remains undetermined because the speed, branch variable, and slope evolve together. Consequently, the explicit estimate in Theorem~\ref{thm:quantitative-approx-flat-no-go} is used only as a universal upper bound in that regime. The next result shows that for a fixed flatness tolerance, the span bounds remain uniform as the outer endpoint increases, excluding an approximately flat trace extending to arbitrarily large radius under the stated hypotheses.
\begin{corollary}\label{cor:no-semi-infinite-approx-flat-trace}
For fixed $R_1>0$ and $0\leqslant\varepsilon<1$, the hypotheses of Theorem~\ref{thm:quantitative-approx-flat-no-go} cannot hold on every interval $[R_1,R_2]$ with $R_2>R_1$.
\end{corollary}
\begin{proof}
For $\ellzero\leqslant 0$, the first bound is independent of $R_2$. For $\ellzero>0$, let us fix any $\alpha\in(0,1-\varepsilon)$. Since $\underline v^2\geqslant0$, the second bound is at most $\mathcal B_+(\varepsilon,\alpha,0)$, also independent of $R_2$. In either case, a finite constant bounds $\log(R_2/R_1)$, which diverges as $R_2\to\infty$.
\end{proof}
These conclusions apply to the constant-$\ell$ class under the stated density and rapidity-shape hypotheses, and they concern the speed measured by ZAMOs. However, an observational interpretation requires the additional relation between that local speed and spectroscopic data. Mathematically, the bounds identify a finite radial limitation that persists despite the analytic freedom to realize traces on individual annular neighbourhoods.


\section{Asymptotic decay and vacuum continuation}
\label{sec:asymptotic-constraint}

We now consider whether a locally realized trace can extend along an unbounded asymptotically flat equatorial end. For a regular branch with $\ellzero\neq 0$ and uniformly subluminal positive speed, the exact metric relations force $rv$ to approach a positive constant fixed by the normalization of $H$. This determines the leading inverse-radius decay of the speed. If the density is nonnegative in an open neighbourhood of the equatorial tail, the decay forces $\eta$ to become constant there, producing a locally flat vacuum region. Analytic continuation then propagates this plateau throughout the connected regular inverse sheet containing it. The resulting metric retains a helical global identification, whose consequences for regular-axis completion and causality we examine. To formulate the asymptotic hypotheses, let $(t,x,y,z)$ be asymptotically Cartesian coordinates in an asymptotically nonrotating frame and set $\mathcal R=(x^2+y^2+z^2)^{1/2}$ and $\varrho=(x^2+y^2)^{1/2}$. We assume the metric falloff
\begin{equation}\label{eq:af-cartesian-falloff}
g_{tt}=-1+O(\mathcal R^{-1}),\qquad
g_{ti}=O(\mathcal R^{-1}),\qquad
g_{ij}=\delta_{ij}+O(\mathcal R^{-1}).
\end{equation}
These estimates provide the asymptotic control needed below. They may be supplemented by the usual derivative conditions when formulating the full isolated-body problem \cite{Regge1974AnnPhys,BeigSimon1980GRG,Beig1987AnnPhys}. The following lemma translates the Cartesian falloff assumptions into estimates for the azimuthal metric components and the ZAMO angular velocity, and relates the orbit-area radius to the cylindrical radius at infinity.
\begin{lemma}\label{lem:af-azimuthal-falloff}
Suppose Eq.~\eqref{eq:af-cartesian-falloff} holds with stationary field
$k=\partial_t$ and axial field
$m=\partial_\phi=-y\partial_x+x\partial_y$. Along the equatorial end
\begin{equation}\label{eq:af-cylindrical-components}
g_{t\phi}=O(1),\qquad
g_{\phi\phi}=\varrho^2+O(\varrho),\qquad
\OmZ=O(\varrho^{-2}).
\end{equation}
The orbit-area radius $r$ satisfies
\begin{equation}\label{eq:canonical-cylindrical-asymptotics}
r^2=\varrho^2+O(\varrho),\qquad
r/\varrho=1+O(\varrho^{-1}).
\end{equation}
The same estimates hold on fixed cones separated from the axis.
\end{lemma}
\begin{proof}
The identities $g_{t\phi}=-y g_{tx}+x g_{ty}$, and $g_{\phi\phi}=\varrho^2+(g_{ij}-\delta_{ij})m^im^j$ give $g_{t\phi}=O(\varrho/\mathcal R)$ and
$g_{\phi\phi}=\varrho^2+O(\varrho^2/\mathcal R)$.
On the equator $\mathcal R=\varrho$, and the shift estimate follows from
$\OmZ=-g_{t\phi}/g_{\phi\phi}$. The exact identity
$r^2=-g_{tt}g_{\phi\phi}+g_{t\phi}^2$ then gives
Eq.~\eqref{eq:canonical-cylindrical-asymptotics}. On a cone separated
from the axis, $\mathcal R$ and $\varrho$ have comparable size.
\end{proof}
For the constant-$\ellzero$ family, the azimuthal metric component gives
\begin{equation}\label{eq:exact-gphiphi-ratio}
\frac{g_{\phi\phi}}{r^2}=\frac{1-v^2}{C e^{\ellzero rv}},\qquad C>0.
\end{equation}
A positive limiting speed would drive this ratio to zero when
$\ellzero>0$ and to infinity when $\ellzero<0$. Asymptotic flatness
requires it to tend to one. The next result gives the full consequence
under uniform subluminality.
\begin{theorem}\label{cor:exact-asymptotic-constraint}
Let a regular exact constant-$\ellzero$ branch with $\ellzero\neq 0$
extend along an asymptotically flat equatorial end. Suppose that for some
$\delta\in(0,1)$ and all sufficiently large $r$ $0<v(r,0)\leqslant1-\delta$ holds. Then, along the equatorial end
\begin{equation}\label{eq:asymptotic-locking-q}
\lim_{r\to\infty}\eta(r,0)=\lim_{r\to\infty}rv(r,0)=q,
\qquad
q=-\frac{\log C}{\ellzero}>0.
\end{equation}
The angular velocity satisfies $\Omega(q)=0$ at this limiting value of $\eta$. Moreover, 
\begin{equation}\label{eq:exact-inverse-radius-decay}
v=O(r^{-1}),\qquad \Omega=O(r^{-2}).
\end{equation}
If also $g_{\phi\phi}/r^2=1+O(r^{-2})$, then
\begin{equation}\label{eq:asymptotic-locking-orders}
\eta=q+O(r^{-2}),\qquad
v=\frac qr+O(r^{-3}),\qquad H=-1+O(r^{-2}).
\end{equation}
In particular, such an end requires $C\neq 1$.
\end{theorem}
\begin{proof}
Lemma~\ref{lem:af-azimuthal-falloff} implies $g_{\phi\phi}/r^2\to1$. Taking the logarithm of Eq.~\eqref{eq:exact-gphiphi-ratio} gives
\begin{equation}\label{eq:log-exact-asymptotic-identity}
\ellzero\eta=\log(1-v^2)-\log C-\log(g_{\phi\phi}/r^2).
\end{equation}
The bound $\delta(2-\delta)\leqslant1-v^2<1$ makes the right-hand side
bounded. Hence, $\eta=O(1)$ and $v=\eta/r\to 0$. Equation
\eqref{eq:log-exact-asymptotic-identity} now gives $\eta\to q=-(\log C)/\ellzero$, with $q\geqslant 0$. The exact Killing-block relation also yields
\begin{equation}\label{eq:omega-from-killing-block}
\Omega=\frac{\eta-g_{t\phi}}{g_{\phi\phi}}.
\end{equation}
Its numerator is bounded and its denominator has order $r^2$, so
$\Omega=O(r^{-2})$. If $q=0$, then $\eta\to 0^+$. Let us recall that integrating $\Omega_{,\eta}=H_{,\eta}/(2\eta)$ with $H=-Ce^{\ellzero\eta}$ gives  \eqref{eq:constant-l-Omega}. This result and the expansion
$\operatorname{Ei}(x)=\gamma_{\rm E}+\log|x|+O(x)$ imply
\begin{equation}\label{eq:omega-log-divergence}
|\Omega|=\frac{C|\ellzero|}{2}|\log\eta|\bigl(1+o(1)\bigr)
\longrightarrow\infty.
\end{equation}
This contradicts the decay just obtained. Therefore, $q>0$, and continuity
of $\Omega(\eta)$ near $q$ gives $\Omega(q)=0$. Under the additional falloff hypothesis, the difference between the
right-hand side of Eq.~\eqref{eq:log-exact-asymptotic-identity} and
$-\log C$ is $O(v^2)+O(r^{-2})=O(r^{-2})$. Thus, $\eta=q+O(r^{-2})$ and $v=q/r+O(r^{-3})$. Expanding $H=-Ce^{\ellzero\eta}$ and using $Ce^{\ellzero q}=1$ gives the stated estimate for $H$. Finally, $C=1$ would give $q=0$.
\end{proof}
This proof uses the exact metric. The logarithmic low-velocity expressions
of \cite{Astesiano2022EPJC} have an expansion parameter $\ellzero\eta$, which grows along a fixed nonzero-speed trace. Their extrapolation to arbitrarily large radius therefore requires a separate remainder estimate. The supplementary material gives the corresponding conditional asymptotic test and explains this loss of uniformity. Combining the asymptotic decay just established with nonnegative density near the equatorial tail forces a vacuum plateau, which extends throughout the connected regular inverse sheet by analytic continuation.
\begin{theorem}\label{thm:eventual-helical-plateau}
Under the hypotheses of Theorem~\ref{cor:exact-asymptotic-constraint}, suppose there exists $R_0>0$ such that every point $(r,0)$ with $r>R_0$ has an open neighbourhood on which $\rho\geqslant 0$. Then, an open neighbourhood of an
equatorial tail satisfies $\eta=q$, and $\rho=0$. If this neighbourhood belongs to a connected regular inverse sheet, then
$\eta=q$ throughout that sheet. After asymptotic normalization, the metric on its connected vacuum plateau is
\begin{equation}\label{eq:helical-plateau-metric}
ds^2=-(dt-q\,d\phi)^2+r^2d\phi^2+dr^2+dz^2
\end{equation}
with $\phi\sim\phi+2\pi$ denoting the identification of angular coordinates differing by $2\pi$, so that $\phi$ parametrizes the closed orbits of the axial symmetry.
\end{theorem}
\begin{proof}
Theorem~\ref{cor:exact-asymptotic-constraint} gives $v=O(r^{-1})$, so both endpoints of the density interval \eqref{eq:density-domain} tend to zero along the equatorial end, whereas $|r\ellzero|\to\infty$. The algebraic density factor is therefore strictly negative at every sufficiently large equatorial point. By continuity, it stays negative on a small connected neighbourhood of each point. Intersecting with the neighbourhood on which $\rho\geqslant 0$ and using the density factorization gives $\nabla\eta=0$ there. The union contains the entire
tail, so the component containing that tail is connected. Moreover, the value of $\eta$ on it equals its asymptotic limit $q$ while the density vanishes. By Theorem~\ref{thm:automatic-analyticity}, the equality $\eta=q$ extends throughout the connected regular inverse sheet containing the equatorial tail. Furthermore, the functions $H$ and $\Omega$ are constant when $\eta$ is constant, and
the reconstruction quadratures make $\mu$ constant. Their values are
fixed by the asymptotic normalization to $H=-1$, $\Omega=0$, and
$\mu=0$. The metric components then become
$g_{tt}=-1$, $g_{t\phi}=q$, and $g_{\phi\phi}=r^2-q^2$, proving
Eq.~\eqref{eq:helical-plateau-metric}.
\end{proof}
In particular, a connected regular inverse sheet containing positive
density cannot also contain the nonnegative-density asymptotic end of this theorem. A finite dust body must therefore be matched to a different vacuum solution, unless at least one of the hypotheses of the continuation theorem fails. The next result signalizes that the plateau metric also constrains continuation towards the axis, since a regular fixed point of the axial symmetry requires $q=0$, in conflict with the positive value imposed by the asymptotic end.
\begin{corollary}\label{cor:helical-regular-axis-exclusion}
If the plateau in Eq.~\eqref{eq:helical-plateau-metric} extends through the same connected flat stationary and axisymmetric region to a regular fixed point of the axial action, then $q=0$. Consequently, the positive-speed end of Theorem~\ref{cor:exact-asymptotic-constraint} has no such regular-axis completion.
\end{corollary}
\begin{proof}
At a regular axial fixed point, the axial Killing field vanishes, so
$g(m,m)=0$ and the orbit-area radius is zero. Continuing $g(m,m)=r^2-q^2$ to that point therefore gives $q=0$, contradicting the conclusion $q>0$ of Theorem~\ref{cor:exact-asymptotic-constraint}. Hence, the plateau associated with that asymptotic end admits no regular-axis completion of the stated kind.
\end{proof}
To describe the effect of the angular identification on the global geometry, we pass to the universal cover, where $\phi$ ranges over the real line. There, the coordinate $T=t-q\phi$ puts the metric in Minkowskian form. The original periodicity becomes the identification $(T,\phi)\mapsto(T-2\pi q,\phi+2\pi)$, so a complete turn around the axis induces a time translation of $-2\pi q$ on the cover. The causal properties of the resulting quotient also depend on its radial extent. On a region $r>r_0>|q|$, the single-valued coordinate $t$ has timelike gradient because $g^{tt}=-1+q^2/r^2<0$, establishing stable causality. If the domain extends into $r<|q|$, the axial orbits become timelike because $g_{\phi\phi}=r^2-q^2<0$, and their periodicity produces closed timelike curves.


\section{Finite boundaries and shell-free matching}
\label{sec:vacuum-matching}

For a finite dust body, completion requires matching the interior geometry to a vacuum exterior across a timelike boundary. We begin by identifying smooth candidate boundaries within the zero-density set and computing their first and second fundamental forms. These expressions give the full Darmois conditions for matching without a surface layer. When the induced metric is continuous, a jump in the second fundamental form determines the surface stress, making explicit the additional requirements beyond vanishing boundary density. The same matching conditions imply continuity of the orbit-area radius and its normal derivative, thereby specifying the corresponding boundary data that the exterior must realize. We start by observing that at a regular point with $r>0$ and $0<v<1$, the density factorization gives
\begin{equation}\label{eq:zero-density-decomposition}
\{\rho=0\}=\mathcal C\cup\mathcal Z_-\cup\mathcal Z_+,
\end{equation}
where
\begin{equation}
\mathcal C=\{(r,z)\mid\nabla_{rz}\eta=0\},\quad
\mathcal Z_-=\left\{(r,z)\,\middle|\,
r\ellzero=-\frac{2v}{1-v^2}\right\},\qquad
\mathcal Z_+=\left\{(r,z)\,\middle|\,
r\ellzero=\frac{2v}{1+v^2}\right\}.
\end{equation}
The critical set may meet either algebraic branch, and the two algebraic branches are disjoint. Moreover, a smooth embedded curve $\gamma$ in this set generates the symmetry-preserving hypersurface
$\Sigma=\mathbb R_t\times S^1_\phi\times\gamma$. It is timelike in every regular region with $r>0$. Indeed, the Killing block has negative determinant, and the meridional tangent has positive
norm. Notice that singular points and degenerate axial orbits require separate analysis. In what follows, it is convenient to parametrize $\gamma$ by Euclidean meridional arclength $\lambda$, so
$\dot r^2+\dot z^2=1$, and set
\begin{equation}\label{eq:euclidean-tangent-normal}
T_{\rm E}=\dot r\,\partial_r+\dot z\,\partial_z,\qquad
N_{\rm E}=-\dot z\,\partial_r+\dot r\,\partial_z,\qquad
\frac{dT_{\rm E}}{d\lambda}=\kappa_{\rm E}N_{\rm E}.
\end{equation}
For the physical meridional metric $\mathsf q=e^\mu(dr^2+dz^2)$, we have 
\begin{equation}\label{eq:proper-tangent-normal}
d\lambda_{\rm p}=e^{\mu/2}d\lambda,\qquad
e_s=e^{-\mu/2}T_{\rm E},\qquad n=e^{-\mu/2}N_{\rm E}.
\end{equation}
Moreover, we use $K(X,Y)=g(\nabla_Xn,Y)$ and the same normal orientation on both sides of a proposed matching. With this choice of tangent and normal fields, we can express the boundary fundamental forms in terms of the Killing block, the meridional metric, and the curvature of the generating curve.
\begin{lemma}\label{lem:symmetry-reduced-boundary-data}
In the intrinsic basis $(\partial_t,\partial_\phi,e_s)$, with
$A,B\in\{t,\phi\}$, the first and second fundamental forms are
\begin{align}
h_{AB}&=G_{AB},&h_{As}&=0,&h_{ss}&=1,
\label{eq:first-fundamental-form-reduced}\\
K_{AB}&=\tfrac{1}{2}n^\nu\partial_\nu G_{AB},&K_{As}&=0,&
K_{ss}&=e^{-\mu/2}\left[\tfrac{1}{2}N_{\rm E}(\mu)-\kappa_{\rm E}\right].
\label{eq:boundary-second-form}
\end{align}
\end{lemma}

\begin{proof}
Restricting the block metric to $\Sigma$ gives \eqref{eq:first-fundamental-form-reduced}. Since the meridional normal is invariant under both Killing flows, $[n,\partial_A]=0$. Applying the identity $2K=\mathcal L_n g$ on tangent vectors therefore yields $2K_{AB}=n^\nu\partial_\nu G_{AB}$, with the mixed components $K_{As}$ vanishing by block orthogonality. For the remaining component, we put $\omega=\mu/2$. The conformal
connection formula for $\mathsf q=e^{2\omega}\delta$ reads
\begin{equation}
\nabla^{\mathsf q}_XY=\nabla^{\rm E}_XY
+X(\omega)Y+Y(\omega)X-\delta(X,Y)\nabla^{\rm E}\omega.
\end{equation}
Since $\nabla^{\rm E}_{T_{\rm E}}N_{\rm E}=-\kappa_{\rm E}T_{\rm E}$, this gives
\begin{equation}
\nabla^{\mathsf q}_{e_s}n=e^{-2\omega}\bigl(N_{\rm E}(\omega)-\kappa_{\rm E}\bigr)T_{\rm E},
\end{equation}
and its $\mathsf q$ inner product with $e_s$ yields the stated formula.
\end{proof}
For a proposed matching, let $-$ denote the dust interior and $+$ a
stationary, axisymmetric, circular vacuum exterior. Moreover, we fix an identification of their boundaries that preserves the Killing basis and the axial period \cite{Vera2002CQG}. Both unit normals point from the dust toward the vacuum. We also  write $[Q]=Q^+-Q^-$. To compare the induced metrics, we parametrize the interior and exterior boundary curves so that the same value of $u$ labels points identified by the matching. With $\gamma_\pm(u)=(r_\pm(u),z_\pm(u))$, their equality requires
\begin{align}
G^+_{AB}&=G^-_{AB},\label{eq:common-parameter-first-form-G}\\
e^{\mu_+}\bigl(\dot r_+^2+\dot z_+^2\bigr)
&=e^{\mu_-}\bigl(\dot r_-^2+\dot z_-^2\bigr).
\label{eq:common-parameter-first-form-meridional}
\end{align}
Only after these conditions hold does the common induced metric determine a common proper arclength and unit tangent $e_s$. This step ensures that the boundary identification preserves physical length along the generating curve. The next result shows that the boundary fundamental forms determine whether the proposed identification gives a matching without a surface layer, produces a surface stress, or fails to define a continuous metric across the interface.
\begin{proposition}\label{prop:matching-alternatives}
Let a $C^2$ rotating-dust interior and a $C^2$ circular vacuum exterior have smooth timelike symmetry-preserving boundaries in regular regions with $r>0$. Under a fixed identification and common normal orientation, the local alternatives are determined by $[h]$ and $[K]$. If $[h]\neq 0$, the identification fails to give a standard Darmois--Israel matching. If $[h]=0$, a shell-free matching is equivalent to $[K_{ab}]=0$. In the common proper-arclength basis, the complete conditions are Eqs.~\eqref{eq:common-parameter-first-form-G}--\eqref{eq:common-parameter-first-form-meridional}
together with
\begin{align}
[n^\nu\partial_{\nu}G_{AB}]&=0,\label{eq:darmois-normal-G}\\
\left[e^{-\mu/2}\left(\tfrac12N_{\rm E}(\mu)-\kappa_{\rm E}\right)\right]
&=0.\label{eq:darmois-meridional-curvature}
\end{align}
If $[h]=0$ and $[K_{ab}]\neq 0$, the continuous gluing has surface stress
\begin{equation}\label{eq:israel-surface-tensor}
S_{ab}=-\frac1{8\pi G_{\rm N}}
\left([K_{ab}]-h_{ab}[K]\right),\qquad [K]=h^{ab}[K_{ab}].
\end{equation}
A smooth component of ${\rho=0}$ provides a candidate boundary whose matching properties are determined by the jumps in the two fundamental forms. Once the induced metrics agree, the surface stress vanishes precisely when $[K_{ab}]=0$.
\end{proposition}

\begin{proof}
Equality of the first fundamental forms is the continuity requirement
for the standard junction construction
\cite{Darmois1927,Israel1966NCB,MarsSenovilla1993CQG}.
After it holds, shell-free matching requires equality of the second
fundamental forms. Lemma~\ref{lem:symmetry-reduced-boundary-data} gives exactly the two derivative conditions stated above, and its mixed components vanish on both sides. When the second fundamental forms differ, the distributional Einstein equation gives Eq.~\eqref{eq:israel-surface-tensor}. Taking its trace on
the three-dimensional boundary yields
\begin{equation}
h^{ab}S_{ab}=\frac{[K]}{4\pi G_{\rm N}}.
\end{equation}
Thus, $S_{ab}=0$ first forces $[K]=0$ and then $[K_{ab}]=0$. Conversely, $[K_{ab}]=0$ implies $[K]=0$, and substitution into \eqref{eq:israel-surface-tensor} gives $S_{ab}=0$. Since the dust velocity is tangent to every symmetry-preserving boundary, i.e.
\begin{equation}\label{eq:dust-tangent-to-boundary}
u=(-H)^{-1/2}(\partial_t+\Omega\partial_\phi),
\end{equation}
it follows that $u\cdot n=0$, and $T_{\mu\nu}n^\nu=0$ regardless of the boundary value of $\rho$. The Gauss and Codazzi equations express the normal Einstein projections through $(h,K)$, so
these projections already agree at a Darmois boundary. Requiring
$\rho=0$ also sets the full matter tensor to zero there, and it leaves the exterior metric and its normal derivatives to be determined by the vacuum field equations and the matching conditions.
\end{proof}
The significance of a zero-density boundary depends on whether a vacuum neighbourhood has already been constructed. Within an existing $C^2$ vacuum collar, both sides of a cut inherit their geometry from the same smooth metric, so the fundamental forms agree automatically. For an isolated smooth component of $\mathcal Z_\pm$ or $\mathcal C$, the data computed above prescribe conditions that a vacuum exterior must satisfy, and the existence of such an exterior remains to be established. Moreover, the matching conditions also require the orbit-area radius and its normal derivative to agree across the boundary, as the following corollary shows.
\begin{corollary}\label{cor:orbit-area-matching}
Let $\varpi=\sqrt{-\det G}$ in the common stationary and axial Killing basis. Every shell-free matching satisfies $[\varpi]=0$, and $[n^\nu\partial_\nu\varpi]=0$. For the present dust interior, $\varpi^-=r$. Hence, an exterior in canonical Weyl coordinates must realize
\begin{equation}\label{eq:weyl-radius-boundary-data}
\varpi^+|_\Sigma=r|_\Sigma,\qquad
n^\nu\partial_\nu\varpi^+|_\Sigma=n^r_\Sigma.
\end{equation}
\end{corollary}
\begin{proof}
Continuity of the first fundamental form gives $[G]=0$, so the determinants agree across $\Sigma$ and hence $[\varpi]=0$. Continuity of the second fundamental form also gives $[n^\nu\partial_\nu G]=0$, since $2K_{AB}=n^\nu\partial_\nu G_{AB}$. To relate these data to the normal derivative of $\varpi$, we apply Jacobi's formula on each side. This leads to
\begin{equation}\label{eq:normal-determinant-identity}
2n^\nu\partial_\nu\log\varpi=\operatorname{tr}!\left(G^{-1}n^\nu\partial_\nu G\right).
\end{equation}
Both $G$ and its normal derivative agree across the boundary, so the right-hand side has the same value on either side. Since $\varpi$ is positive and continuous there, it follows that $[n^\nu\partial_\nu\varpi]=0$. The stated boundary data now follow from $\varpi^-=r$ and $n^\nu\partial_\nu r=n^r$ in the dust interior, with $n^r_\Sigma$ denoting this interior normal component evaluated on $\Sigma$.
\end{proof}
The supplementary material presents the surface tensor in a symmetry-adapted basis and gives explicit formulas for a cylindrical boundary. To match the dust interior to a vacuum exterior without a surface layer, the two sides must have the same induced metric and satisfy the normal matching conditions in Eqs.~\eqref{eq:darmois-normal-G}--\eqref{eq:darmois-meridional-curvature}.


\section{Topology and compatibility of the exterior}
\label{sec:topology-exterior}

The matching data derived above must be realized by a vacuum solution on the entire exterior domain. We first examine how the topology of that domain affects reconstruction. Under the stated assumptions, a body with a disk-shaped meridional section separated from the axis is a solid torus, and the resulting exterior topology imposes period conditions for single-valued potentials. These conditions enter our definition of an admissible isolated vacuum completion, together with the requirements of boundary matching, geometric regularity, and asymptotic flatness. We then examine the compatibility between the boundary values and normal derivatives supplied by the interior. Linearizing the Ernst equations about Minkowski space on a fixed smooth bounded domain away from the axis gives a precise characterization of admissible boundary data through the Dirichlet-to-Neumann map. This provides a compatibility test for the linearized problem. Constructing a nonlinear exterior satisfying all the conditions for an isolated completion remains an open existence problem.


\subsection{Toroidal bodies and period conditions}

The following theorem relates the toroidal topology of a body separated from the axis to the period conditions required for single-valued exterior potentials.
\begin{theorem}\label{thm:toroidal-body-periods}
Let $\mathscr S$ be a smooth spacelike hypersurface invariant under an effective axial $U(1)$ action, with a regular axis and no exceptional orbits in the region considered. Let $Q=\mathscr S/U(1)$ be the orbit surface and $\pi$ the quotient projection. Suppose $Q$ is diffeomorphic to a half-plane and a compact connected invariant body $\Omega_{\rm mat}\subset\mathscr S$ projects under $\pi$ onto a smoothly embedded closed disk $B$ contained in the interior of $Q$. Then $\Omega_{\rm mat}\simeq D^2\times S^1$, $\partial\Omega_{\rm mat}\simeq T^2$, and the exterior orbit surface $Q^+=Q\setminus\operatorname{int}B$
satisfies $\pi_1(Q^+)\simeq\mathbb Z$, and $H^1_{\rm dR}(Q^+)\simeq\mathbb R$. Assume that the exterior metric satisfies the vacuum Einstein equations, is circular, and has a Lorentzian Killing block. Let $\mathsf q$ be its orbit metric and $\varpi=\sqrt{-\det G}$. The one-forms $*_{\mathsf q}d\varpi$ and the descended stationary twist form $\omega_k$ are closed. If $\gamma$ generates the first homology of $Q^+$, global potentials satisfying $d\zeta=*_{\mathsf q}d\varpi$ and $d\chi=\omega_k$ exist exactly when
\begin{equation}\label{eq:topological-period-tests}
P_\zeta=\oint_\gamma *_{\mathsf q}d\varpi=0,\qquad
P_\chi=\oint_\gamma\omega_k=0.
\end{equation}
At points with $d\varpi\neq 0$, $(\varpi,\zeta)$ are local Weyl
coordinates. For $(\varpi,\zeta)$ to define coordinates on the entire exterior orbit surface, the coordinate map must also be injective, so that distinct points have distinct coordinate pairs.
\end{theorem}

\begin{proof}
Recall that $B=\pi(\Omega_{\rm mat})$ is the closed disk representing the body in the orbit surface. Since $B$ lies in the interior of $Q$, it is separated from the axis. The absence of exceptional orbits therefore makes the axial action free over $B$, so $\pi^{-1}(B)$ is a principal circle bundle over this disk. Contractibility of $B$ makes the bundle trivial, giving
$\Omega_{\rm mat}\simeq D^2\times S^1$ and
$\partial\Omega_{\rm mat}\simeq T^2$. The exterior orbit surface $Q^+=Q\setminus\operatorname{int}B$ is a half-plane with one interior disk removed and hence has the homotopy type of a circle. Consequently, $\pi_1(Q^+)\simeq\mathbb Z$ and
$H^1_{\rm dR}(Q^+)\simeq\mathbb R$. This topology determines whether the closed forms supplied by the vacuum equations admit global potentials. For a circular vacuum metric, $\varpi$ is harmonic with respect to $\mathsf q$, so $d(**{\mathsf q}d\varpi)=0$.
For the stationary Killing field, define
\begin{equation}\label{eq:twist-one-form}
(\omega_k)*\tau=\epsilon_{\tau\nu\rho\sigma}k^\nu\nabla^\rho k^\sigma.
\end{equation}
The Killing identity implies $d\omega_k=0$ in vacuum. Moreover, $\omega_k$ is invariant under both Killing flows and, by circularity, annihilates both Killing fields. These properties ensure that $\omega_k$ defines a closed one-form on the orbit surface $Q^+$ \cite{Wald1984GR,Stephani2003Exact}. Let us recall that a closed one-form admits a global potential precisely when all its periods vanish. Since $[\gamma]$ generates the first homology of $Q^+$, this condition reduces to vanishing of the integral along $\gamma$. Applying this criterion to $*_{\mathsf q}d\varpi$ and $\omega_k$ gives the necessary and sufficient conditions \eqref{eq:topological-period-tests} for the existence of $\zeta$ and $\chi$. Where $d\varpi\neq 0$, the relation $d\zeta=*_{\mathsf q}d\varpi$ makes $d\varpi$ and $d\zeta$ linearly independent, so $(\varpi,\zeta)$ defines local coordinates. For these coordinates to form a global chart, $d\varpi$ must remain nonzero throughout the domain and the coordinate map must be injective, assigning distinct coordinate pairs to distinct points.
\end{proof}
For $n$ disjoint disk-like meridional bodies, each removed disk contributes an independent cycle to the exterior orbit surface. A closed reconstruction form therefore admits a global potential precisely when its integral vanishes around each of these $n$ cycles. The exterior problem must be formulated on this multiply connected domain so that all the required period conditions are included. These constraints also enter the reconstruction of the meridional conformal factor. Writing its local vacuum quadratures as $d\mu=\alpha_\mu$, their compatibility gives $d\alpha_\mu=0$. The existence of a single-valued $\mu$ then depends on the periods of this closed form. For the one-body exterior, the condition reduces to
\begin{equation}\label{eq:meridional-period-test}
P_\mu=\oint_\gamma\alpha_\mu=0.
\end{equation}
The exterior metric obtained through this reconstruction must then be matched to the dust metric determined by the auxiliary potential $\mathcal F$, with equality of the first and second fundamental forms at their common boundary.


\subsection{The global completion problem}
\label{subsec:global-vacuum-compatibility}

We now specify what it means to complete the dust interior to an isolated spacetime. This requires a vacuum exterior that matches the interior without a surface layer and extends to an asymptotically flat end with the prescribed regularity and topology. The following definition collects the boundary conditions and the global requirements on that exterior.
\begin{definition}\label{def:admissible-vacuum-completion}
An admissible isolated vacuum completion across $\Sigma$ is a connected stationary, axisymmetric, circular exterior $(\mathcal M^+,g^+)$ with commuting Killing fields $k^+$ and $m^+$ satisfying the following conditions.
\begin{enumerate}
\item 
The metric is $C^2$ up to $\Sigma$, smooth in the exterior, and
satisfies $R_{\mu\nu}[g^+]=0$. Axial orbits have period $2\pi$.
\item 
The boundary is smooth, connected, timelike, and symmetry-preserving.
Its identification with the dust boundary preserves the Killing basis and axial period, and the first and second fundamental forms agree.
\item 
The exterior has no additional boundary, horizon, curvature
singularity, or conical defect. Every axis segment in its closure is
regular.
\item 
The Weyl, twist, and meridional reconstruction forms have vanishing
periods, so the corresponding potentials are single-valued. Wherever a global Weyl chart is used, its coordinate map is noncritical and
injective. If the coordinate map is additionally required to be proper, the inverse image of every compact subset of the specified coordinate range must be compact.
\item 
There is one asymptotically flat end satisfying
Eq.~\eqref{eq:af-cartesian-falloff}. Moreover, the stationary field is normalized to unit time translation and the angular coordinate is asymptotically nonrotating.
\end{enumerate}
\end{definition}
The exact density cutoff identifies a candidate boundary within the interior solution. Completing the spacetime across this boundary without a surface layer requires a vacuum metric that realizes the full matching data
\begin{equation}\label{eq:full-matching-data-summary}
\left(G_{AB},\ n^\nu\partial_\nu G_{AB},\ K_{ss}\right)\big|_\Sigma
\end{equation}
together with the common meridional line element. In the Ernst formulation, these conditions prescribe both boundary values and normal derivatives of the potentials. Their simultaneous prescription therefore raises a compatibility question, since the vacuum field equations relate the two. We examine this relation at the linearized level on a fixed smooth bounded orbit domain $D$ with $\overline D\subset{r>0}$. After fixing the stationary and axial gauge and the additive constant of the twist potential, we describe the vacuum field by the Ernst variables $X=(f,\chi)$, where
$f=-g(k,k)>0$. The equations governing these potentials follow from the action
\begin{equation}\label{eq:ernst-sigma-action}
I_{\rm E}[f,\chi]=\frac12\int_D\frac r{f^2}
\left(|\nabla f|^2+|\nabla\chi|^2\right)\,dr\,dz.
\end{equation}
Linearizing about the Minkowski solution $(f,\chi)=(1,0)$ gives a Dirichlet-to-Neumann map relating boundary values to conormal derivatives. This map characterizes which pairs of linearized boundary data can be realized by a solution on $D$. The following theorem identifies the compatible boundary data as the graph of this Dirichlet-to-Neumann map and establishes its self-adjointness and the Lagrangian structure of the graph.
\begin{theorem}\label{thm:linearized-ernst-calderon}
At the Minkowski point $(f,\chi)=(1,0)$, each Jacobi field
$u\in\{\delta f,\delta\chi\}$ satisfies
\begin{equation}\label{eq:linearized-ernst-equation}
\partial_r(r\partial_ru)+\partial_z(r\partial_zu)=0.
\end{equation}
Let $D$ be smooth and bounded with $\overline D\subset\{r>0\}$.
Its Dirichlet-to-Neumann map $\Lambda_{\rm E}$ sends the boundary value to the outward conormal datum $r\partial_nu$. For the pair of Jacobi fields, it is self-adjoint with principal symbol
\begin{equation}
\sigma_1(\Lambda_{\rm E})(x,\xi)=r(x)|\xi|\operatorname{Id}_{\mathbb R^2}.
\end{equation}
Its graph coincides with the space of all Cauchy data of solutions on $D$ and forms an exact Lagrangian subspace of $H^{1/2}(\partial D,\mathbb R^2)\oplus H^{-1/2}(\partial D,\mathbb R^2)$ with respect to the canonical boundary symplectic form. Hence, prescribed linearized Ernst boundary values and conormal data
extend over $D$ precisely when they belong to this graph.
\end{theorem}
\begin{proof}
The quadratic part of Eq.~\eqref{eq:ernst-sigma-action} is
\begin{equation}
I_{\rm E}^{(2)}[u_1,u_2]
=\frac12\int_D r\bigl(|\nabla u_1|^2+|\nabla u_2|^2\bigr)\,dr\,dz.
\end{equation}
Variation gives Eq.~\eqref{eq:linearized-ernst-equation} for each
component. Since $r$ is positive and bounded away from zero on
$\overline D$, the Dirichlet problem is uniformly elliptic and uniquely solvable. Uniqueness also follows directly by integrating the equation against a solution with zero boundary value, which gives
$\int_Dr|\nabla u|^2=0$. For solutions $u$ and $w$, Green's identity gives
\begin{equation}
\langle u|_{\partial D},\Lambda_{\rm E}w|_{\partial D}\rangle
=\int_Dr\nabla u\cdot\nabla w=\langle\Lambda_{\rm E}u|_{\partial D},w|_{\partial D}\rangle.
\end{equation}
Thus, $\Lambda_{\rm E}$ is self-adjoint. For the canonical symplectic
form $\sigma((\phi,p),(\psi,q))=\langle\phi,q\rangle-\langle\psi,p\rangle$, the graph is isotropic because the symplectic form vanishes on every pair of elements of the graph. If $(\phi,p)$ is symplectically orthogonal to the graph, then
$\langle\phi,\Lambda_{\rm E}\psi\rangle=\langle\psi,p\rangle$
for all $\psi$. Self-adjointness gives $p=\Lambda_{\rm E}\phi$.
The graph therefore equals its symplectic orthogonal and is Lagrangian. Moreover, the quadratic functional
$\langle\phi,\Lambda_{\rm E}\phi\rangle/2$ generates the graph,
proving exactness. Freezing $r$ and flattening the boundary reduces the principal part to the half-space Laplace equation. Its decaying Fourier mode has outward normal derivative $|\xi|$ times its boundary value. Multiplication by the frozen coefficient gives the symbol $r|\xi|$ for each of the two components.
\end{proof}
The theorem provides a compatibility criterion for the linearized Ernst variables on a fixed bounded domain. Applying it to data from a dust interior requires a boundary identification that fixes the relative orientation of the normals, so that interior derivatives can be compared consistently with the exterior Dirichlet-to-Neumann map. To formulate a complete isolated exterior problem, one must also prescribe the behaviour at infinity, the axis and topology, and the independent orbit-area and meridional metric data. The supplementary material sets out the additional hypotheses needed to formulate the nonlinear matching problem. Earlier work addresses these questions under several choices of boundary data and exterior geometry. For non-rigid dust, Zingg, Aste, and Trautmann analyzed the overdetermination arising from the simultaneous prescription of metric and normal data \cite{ZinggAsteTrautmann2007ASTP}. One approach to exterior construction is to prescribe Dirichlet data, for which existence and uniqueness results are available under restricted hypotheses \cite{SchaudtPfister1996PRL,Schaudt1998CMP,Klenk1999ANZIAM}. Explicit constructions can also exploit the structure of particular interiors, as in Marvan's Papapetrou-vacuum matchings for rigid van Stockum dust with suitable boundaries \cite{Marvan2023JGP}. Once an isolated exterior exists, the theorem of Mars and Senovilla establishes uniqueness with the matching data and twist normalization fixed \cite{MarsSenovilla1998MPLA}. The asymptotic assumptions remain essential to this distinction between construction and uniqueness, since procedures for prescribed interiors may permit other outer geometries \cite{HauserErnst2001GRG}. A further approach treats compatibility perturbatively, with necessary and sufficient conditions established through second order about static isolated bodies \cite{MacCallumMarsVera2007PRD}. These results provide precedents for the exterior analysis, while the nonlinear existence problem for the zero-density boundaries considered here remains open. For the present family, the local realization theorems provide regular positive-density neighbourhoods of prescribed traces, whose radial spans are constrained under the stated shape hypotheses. The asymptotic analysis determines a limit to their continuation. Extending the same connected regular sheet with nonzero $\ellzero$ to an asymptotically flat end under the density and subluminality assumptions would force a vacuum plateau throughout that sheet. Completion of a finite dust body therefore requires a separate vacuum exterior realizing the full Darmois data, together with the topology, period, regularity, and asymptotic conditions specified above.


\section{From local rotation profiles to spectroscopic observables}\label{sec:observable-scope}

To relate a local speed profile to spectroscopic measurements, we must specify the emitter, the observer, and the null geodesic connecting them. Let $x_{\rm e}$ and $x_{\rm o}$ denote the emission and reception events, and let $p$ be the tangent to the connecting future-directed, affinely parametrized null ray. The stationary and axial Killing fields give the conserved quantities $E=-g(p,k)$, and  $L=g(p,m)$. For an emitter and observer in circular motion, we write their four-velocities as
\begin{equation}\label{eq:circular-source-observer}
u_a=U_a(k+\Omega_a m),\qquad
U_a=\left[-g(k+\Omega_a m,k+\Omega_a m)|_{x_a}\right]^{-1/2},
\quad a\in\{{\mathrm e,\mathrm o}\}.
\end{equation}
The conserved quantities then allow us to express the measured frequency ratio in terms of these motions and the ray's impact parameter.
\begin{proposition}\label{prop:invariant-redshift-formula}
For $E\neq 0$, define $b=L/E$. The frequencies measured at emission and reception satisfy
\begin{equation}\label{eq:exact-killing-redshift}
1+z_{\rm obs}=\frac{\omega_{\rm e}}{\omega_{\rm o}}
=\frac{U_{\rm e}}{U_{\rm o}}
\frac{1-\Omega_{\rm e}b}{1-\Omega_{\rm o}b}.
\end{equation}
For a dust emitter, $U_{\rm e}=(-H_{\rm e})^{-1/2}$.
\end{proposition}
\begin{proof}
Along the ray, the geodesic and Killing equations give
$\nabla_p[g(p,k)]=p^\mu p^\nu\nabla_\mu k_\nu=0$,
with the analogous relation for $m$. Thus, $E$ and $L$ have the same values at both endpoints, where the measured frequencies are
$\omega_a=-g(p,u_a)=U_a(E-\Omega_a L)$. Substituting $L=bE$ and cancelling the common factor $E$ gives the stated ratio. For a dust emitter, the normalization follows from
$g(k+\Omega_{\rm e}m,k+\Omega_{\rm e}m)=H_{\rm e}$.
\end{proof}
Consider now a fixed ray reaching the asymptotically flat end of a completed spacetime. For observers approaching rest in the normalized asymptotic frame, $U_{\rm o}\to 1$. If also $\Omega_{\rm o}b\to 0$, the frequency ratio for a dust emitter reduces in this limit to
\begin{equation}\label{eq:dust-redshift-asymptotic-observer}
1+z_{\rm obs}\longrightarrow
\frac{1-\Omega_{\rm e}b}{\sqrt{-H_{\rm e}}}.
\end{equation}
Evaluating this expression requires the impact parameter of the ray connecting the emitter to the observer. Determining that ray is a null-geodesic boundary-value problem involving the completed metric and the viewing geometry. The dependence on the ray direction is already present in the local frequency measurement. At emission, let $u_{\rm Z}$ be the ZAMO four-velocity and $e_{\hat\phi}$ its unit azimuthal vector. Let $u_{\rm e}=\Gamma(u_{\rm Z}+ve_{\hat\phi})$ and
$p=\omega_{\rm Z}(u_{\rm Z}+\widehat\nu)$,
where $\Gamma=(1-v^2)^{-1/2}$, $\omega_{\rm Z}$ is the frequency measured by the ZAMO, and $\widehat\nu$ is the unit spatial propagation direction in that frame. Contracting these expressions gives
\begin{equation}\label{eq:local-doppler-direction-factor}
\omega_{\rm e}=\Gamma\omega_{\rm Z}(1-v\cos\alpha),\qquad
\cos\alpha=g(\widehat\nu,e_{\hat\phi}).
\end{equation}
Together with the frequency-ratio formula, this relation shows how the emission direction and observer motion enter the spectroscopic prediction alongside the local speed. The flatness assumptions in the span theorems constrain the equatorial ZAMO-speed trace. Relating that trace to a measured rotation curve therefore requires specifying these additional data and propagating the rays through a completed spacetime. The supplementary material gives a detailed formulation of this prediction problem.


\section{Conclusions and Outlook}\label{sec:conclusions_outlook}

The constant-$\ell$ dust family provides a setting in which local realizability and global completion can be studied within the exact field equations. The linearizing potential reduces the velocity equation to a linear elliptic problem, while the density inequality determines the admissible inverse sheet and identifies the fold at its lower boundary. On regular annular neighbourhoods, this structure permits analytic realizations of exactly flat equatorial traces with positive density. Fixing the reconstruction constants then determines the corresponding local dust metric. Automatic analyticity controls the propagation of these solutions, and the instability of auxiliary elliptic Cauchy continuation quantifies their sensitivity to the prescribed data. Extending these local constructions introduces further geometric restrictions. Under the stated axis hypotheses, metric regularity forces $\ellzero=0$, and analyticity excludes a positive exactly flat trace on a connected equatorial branch reaching the axis. For annular traces, nonnegative density and invariant transverse concavity yield finite radial-span bounds at any fixed logarithmic-slope tolerance below unity. The bound for the nonpositive class is attained on the $\ellzero=0$ branch, while the positive-$\ellzero$ exactly flat supremum is approached by analytic local solutions. These results quantify the extent to which nearly flat profiles can persist under the specified density and shape assumptions. The asymptotic analysis gives a further restriction on continuation. For a regular branch with nonzero $\ellzero$ and uniformly subluminal positive speed, a standard asymptotically flat equatorial end forces $rv$ to approach a positive constant. Nonnegative density in an open neighbourhood of the tail then produces a locally flat vacuum plateau, which extends throughout the connected regular inverse sheet. Its helical identification retains global geometric information and restricts both its causal domain and its possible continuation to a regular axis. For a finite body, completion requires a vacuum exterior realizing the full Darmois data. The zero-density set supplies candidate boundaries, whose first and second fundamental forms determine the matching requirements. The topology of a body separated from the axis introduces additional period conditions for single-valued exterior potentials. Within this framework, the linearized Ernst analysis characterizes compatible boundary values and conormal derivatives about Minkowski space on a fixed smooth bounded domain away from the axis. It also provides a precise starting point for studying the compatibility of interior data with the vacuum equations.

The principal open problem is to construct a finite dust body together with an admissible isolated vacuum exterior. This requires extending the local annular solutions to a region with a smooth closed material boundary and finding an exterior that realizes its induced metric and extrinsic curvature. A natural sequel would formulate this as a coupled problem in which the interior data, boundary shape, and exterior geometry are determined together. The first step would be to identify which variations of the interior and boundary produce compatible exterior data, followed by an analysis of whether that compatibility persists for the nonlinear Ernst equations. The period conditions, axis regularity, and asymptotic normalization would have to be incorporated throughout this construction. A second open question concerns the optimal radial span for positive $\ellzero$ at nonzero flatness tolerance. Resolving it requires controlling the coupled evolution of the speed, its logarithmic slope, and the branch variable, together with constructing analytic solutions that approach the resulting bound. The auxiliary Cauchy instability also motivates the development of reconstruction methods that control sensitivity while preserving density positivity and inverse-sheet regularity. The scope of the present conclusions is the constant-$\ell$ class under the stated hypotheses. Extending the analysis to variable $\ell$ would allow us to investigate whether approximately constant annular behaviour can connect to a regular axis while satisfying $H'(0)=0$. Other matter models and asymptotic geometries require corresponding existence and continuation analyses. Once a completed spacetime is available, its observational predictions can be obtained by specifying emitter and observer motions and solving the connecting null-geodesic problem. This would determine how the locally prescribed rotation profile is reflected in the measured frequency shifts.

\section*{Data availability}

No observational or experimental datasets were used in this theoretical study. The analytical derivations are provided in the manuscript and the Supplementary Material. The accompanying source package includes
the density-boundary and galactic-calibration tables. The verification code is available in the repository at
\begin{center}
  \url{https://github.com/dutykh/Constant-Ell-Rotating-Dust}
\end{center}

\bibliography{DB-DD-ConstantEll-RotatingDust-PartI}
\end{document}
